\documentclass[openany,a4paper,12pt]{book}

% --- Lingua e Font ---
\usepackage[utf8]{inputenc}
\usepackage[T1]{fontenc} 
\usepackage[italian]{babel}
\usepackage{lmodern} % Font più nitido in PDF rispetto al default
\usepackage{microtype} % Migliora spaziature e sillabazione

% --- Colori (Definiti prima per usarli ovunque) ---
\usepackage[x11names, dvipsnames, table]{xcolor}
\definecolor{PoliBlue}{RGB}{0, 45, 114} % Blu accademico elegante
\definecolor{OpticRed}{RGB}{200, 10, 30} % Rosso "laser" per i link
\definecolor{ForestGreen}{RGB}{34, 139, 34}
\definecolor{CodeBack}{RGB}{245, 245, 245}
\definecolor{BoxGray}{RGB}{248, 249, 250}

% --- Matematica e Fisica ---
\usepackage{amsmath, amssymb, amsfonts, amsthm}
\usepackage{mathtools, cancel, bm}

% Non usare physics insieme a siunitx: i due pacchetti definiscono entrambi \qty e generano conflitti.
% Per le unità di misura usa solo siunitx.
\usepackage{siunitx}
\sisetup{
    output-decimal-marker={,},
    separate-uncertainty=true, % Per scrivere (10,5 ± 0,1) m
    per-mode=symbol, % Scrive m/s invece di m s^{-1}
    inter-unit-product = \ensuremath{{}\cdot{}} % Punto tra le unità
}

% --- Formattazione Testo e Layout ---
\usepackage{ragged2e}
\usepackage{quoting}
\usepackage{booktabs} % Per tabelle professionali
\usepackage{enumitem}
\usepackage{array, tabularx}
\usepackage[a4paper, margin=2.5cm, headheight=15pt]{geometry} 
\setlength{\parindent}{0pt}
\setlength{\parskip}{0.75em}
\usepackage{placeins}

% --- Titoli di Sezione (Stile Libro Universitario Moderno) ---
\usepackage{titlesec}
\titleformat{\chapter}[display]
  {\normalfont\bfseries\color{PoliBlue}}{\filleft\Huge\thechapter}{2ex}
  {\titlerule\vspace{2ex}\Huge\filleft}
  [\vspace{2ex}\titlerule]
\titleformat{\section}
  {\normalfont\Large\bfseries\color{PoliBlue}}{\thesection}{1em}{}
\titleformat{\subsection}
  {\normalfont\large\bfseries\color{PoliBlue}}{\thesubsection}{1em}{}

% --- Grafica, TikZ e Plot ---
\usepackage{graphicx}
\usepackage{float} 
\usepackage{circuitikz}
\usepackage{pgfplots}
\pgfplotsset{compat=1.18}
\usepackage{subcaption} 
\usepackage{tikz}
\usepackage{tikz-3dplot}
\usetikzlibrary{positioning, arrows.meta, angles, quotes, trees, patterns, calc, matrix}

\pgfplotsset{
    standard/.style={
        axis lines=middle,
        xlabel={$t$}, ylabel={$y$},
        grid=both,
        grid style={line width=.1pt, draw=gray!20},
        major grid style={line width=.2pt, draw=gray!40},
        samples=100,
        enlargelimits=true,
        label style={font=\small},
        tick label style={font=\tiny}
    }
}

% --- Box ed Evidenziazioni (Tcolorbox) ---
\usepackage[most]{tcolorbox}

% Stile per le Note (si spezzano se a fine pagina)
\newtcolorbox{nota}[1][]{
  enhanced, breakable,
  left=1em, right=1em, top=1ex, bottom=1ex,
  colback=BoxGray, colframe=gray!50,
  borderline west={3pt}{0pt}{gray},
  fontupper=\itshape,
  title={\textbf{Nota}}, coltitle=black, attach title to upper={\quad},
  #1
}

% Stile per le Definizioni
\newtcolorbox{definizione}[1][]{
  enhanced, breakable,
  left=1em, right=1em, top=1ex, bottom=1ex,
  colback=PoliBlue!5, colframe=PoliBlue!80,
  borderline west={3pt}{0pt}{PoliBlue},
  title={\textbf{Definizione}}, coltitle=PoliBlue!80!black, attach title to upper={\quad},
  #1
}

% Stile per i Teoremi/Formule importanti
\newtcbtheorem[no counter]{teorema}{Teorema}%
{enhanced, breakable, colback=OpticRed!5, colframe=OpticRed, fonttitle=\bfseries}{th}

% --- Configurazione LISTINGS (Python per Analisi Dati) ---
\usepackage{listings}
\lstset{
    language=Python,
    backgroundcolor=\color{CodeBack},
    basicstyle=\ttfamily\small,
    keywordstyle=\color{PoliBlue}\bfseries,
    stringstyle=\color{OpticRed},
    commentstyle=\color{ForestGreen}\itshape,
    numberstyle=\tiny\color{gray},
    numbers=left,
    numbersep=10pt,
    frame=leftline,
    framerule=2pt,
    rulecolor=\color{PoliBlue!50},
    showspaces=false,
    showstringspaces=false,
    breaklines=true,
    breakatwhitespace=false,
    columns=fullflexible, % Permette un copia-incolla perfetto dal PDF
    keepspaces=true,
    literate={à}{{\`a}}1 {è}{{\`e}}1 {ì}{{\`i}}1 {ò}{{\`o}}1 {ù}{{\`u}}1 {±}{{$\pm$}}1
}

% --- Intestazioni e Piè di Pagina ---
\usepackage{fancyhdr}
\pagestyle{fancy}
\fancyhf{}
\renewcommand{\headrulewidth}{0.4pt}
\renewcommand{\chaptermark}[1]{\markboth{\thechapter.\ #1}{}}
\renewcommand{\sectionmark}[1]{\markright{\thesection.\ #1}}
\fancyhead[LE]{\small\textbf{\nouppercase{\leftmark}}} % Capitolo a sx (pagine pari)
\fancyhead[RO]{\small\textbf{\nouppercase{\rightmark}}} % Sezione a dx (pagine dispari)
\fancyfoot[C]{\small\textbf{\thepage}}

\fancypagestyle{plain}{%
    \fancyhf{}%
    \renewcommand{\headrulewidth}{0pt}%
    \fancyfoot[C]{\small\textbf{\thepage}}%
}

% --- Link e Indici (DEVE ESSERE L'ULTIMO PACCHETTO) ---
\usepackage{hyperref}
\usepackage[italian]{cleveref}
\usepackage{bookmark}
\hypersetup{
    colorlinks=true,
    linkcolor=PoliBlue, 
    urlcolor=OpticRed,
    citecolor=ForestGreen,
    pdftitle={Appunti Fondamenti Fisici di Tecnologia Moderna},
    pdfauthor={Agata Trovato, Federico Dogo}
}

% ==================================================
% INIZIO DOCUMENTO
% ==================================================
\begin{document}

\begin{titlepage}
    % --- Grafica di Sfondo (Richiede TikZ) ---
    \begin{tikzpicture}[remember picture, overlay]
        % Barra laterale blu
        \fill[PoliBlue] (current page.north west) rectangle ([xshift=1.5cm]current page.south west);
        
        % Badge Anno Accademico in alto a destra
        \node[fill=OpticRed, text=white, font=\sffamily\bfseries\Large, anchor=north east, minimum width=4cm, minimum height=1.5cm] 
            at ([yshift=-3cm]current page.north east) {2026 / 2027};
            
        % Onde di segnale decorative in basso (tema del corso: segnali/frequenze)
        \begin{scope}[shift={(current page.south west)}, xshift=1.5cm]
            % Onda 1 (Rossa, lenta)
            \draw[OpticRed, line width=2pt, opacity=0.6] 
                (0, 4) sin (4, 5.5) cos (8, 4) sin (12, 2.5) cos (16, 4) sin (20, 5.5);
            % Onda 2 (Blu, media)
            \draw[PoliBlue, line width=1.5pt, opacity=0.4] 
                (0, 3) sin (2.5, 1.5) cos (5, 3) sin (7.5, 4.5) cos (10, 3) sin (12.5, 1.5) cos (15, 3) sin (17.5, 4.5) cos (20, 3);
            % Onda 3 (Verde, veloce - rumore/armoniche)
            \draw[ForestGreen, line width=1pt, opacity=0.5] 
                (0, 4.5) sin (1.5, 3) cos (3, 4.5) sin (4.5, 6) cos (6, 4.5) sin (7.5, 3) cos (9, 4.5) sin (10.5, 6) cos (12, 4.5) sin (13.5, 3) cos (15, 4.5) sin (16.5, 6) cos (18, 4.5) sin (19.5, 3);
        \end{scope}
    \end{tikzpicture}
    
    % --- Testo del Titolo ---
    \vspace*{4cm}
    \begin{flushleft}
        \hspace{1cm} % Sposta il testo a destra per non sovrapporsi alla barra blu
        \begin{minipage}{0.85\textwidth}
            {\fontsize{35}{45}\selectfont \textbf{\textcolor{PoliBlue}{Fondamenti Fisici}}}\\[0.3cm]
            {\fontsize{35}{45}\selectfont \textbf{\textcolor{PoliBlue}{di Tecnologia Moderna}}}\\[0.8cm]
            
            % Linea separatrice
            \textcolor{gray!50}{\rule{\textwidth}{3pt}}\\[1cm]
            
            {\huge \textbf{Michele Tedeschini}}\\[0.5cm]
            {\Large \textcolor{gray!80}{\textit{Appunti del corso}}}
        \end{minipage}
    \end{flushleft}
    
    \vfill
\end{titlepage}

% ==================================================
% INDICE GENERALE
% ==================================================
\cleardoublepage
% Aggiunge l'indice ai segnalibri del PDF a sinistra (fuori dal testo)
\pdfbookmark[0]{Indice}{indice} 
% Genera l'indice cliccabile
\tableofcontents
\cleardoublepage

\chapter{Introduzione ai Segnali}

Un segnale è una funzione che convoglia una informazione riguardo un fenomeno. Tipicamente ciò avviene sottoforma di una onda elettromagnetica o di una corrente elettrica che trasporta dati da un sistema o una rete a un altro terminale.

Consideriamo un segnale "semplice": sinusoidale periodico. Per esempio pensiamo a una persona che vocalizzi una nota per un certo lasso di tempo (qui non tratteremo segnali più complessi). Il segnale viene rappresentato come un onda caratterizzata da lunghezza d'onda $\lambda$, periodo $T$, e ampiezza $A$:

\begin{figure}[H]
    \centering
    \begin{tikzpicture}[xscale=1.5, yscale=1.5]
        % Assi tempo
        \draw[-{Stealth}, thick, PoliBlue!80] (0,-1.5) -- (0,1.5);
        \draw[-{Stealth}, thick, PoliBlue!80] (0,0) -- (3.5,0) node[right, text=PoliBlue] {$t \, , \quad x$};
        
        % Curva Coseno
        \draw[thick, orange!70!yellow, domain=0:3, samples=100] plot (\x, {cos(\x*360/2)});
        
        % Quote (Ampiezza e Periodo)
        \draw[|-|, thin, gray] (-0.2, 0) -- (-0.2, 1) node[midway, left, text=PoliBlue] {$A$};
        \draw[|-|, thin, gray] (0, 1.2) -- (2, 1.2) node[midway, above, text=teal] {$T \, , \quad \lambda$};
        
        % Equazione a destra
        \node[right] at (4, 0.5) {$s = A \cos (\omega t) = A \cos \left(\frac{2\pi}{T}t\right)$};
    \end{tikzpicture}
\end{figure}

\begin{itemize}[label=\textcolor{PoliBlue}{\textbullet}]
    \item Per definizione la frequenza è l'inverso del periodo $\nu = 1 / T$ ed è indipendente dall'ampiezza.
    \item Da ultimo, la potenza $P$ di un segnale periodico $s$ ci è data da $P = \langle s^2(t) \rangle$ (media temporale).
\end{itemize}

Possiamo allora rappresentare il nostro segnale in un piano diverso da quello precedente: al posto del tempo $t$ consideriamo la frequenza $\nu$ e al posto dell'ampiezza $A$ il segnale sarà caratterizzato dalla potenza $P$.

\begin{figure}[H]
    \centering
    \begin{tikzpicture}[xscale=1.2, yscale=1.2]
        % --- DOMINIO DEL TEMPO (Sinistra) ---
        \draw[-{Stealth}, thick, PoliBlue!80] (0,-1.5) -- (0,1.5);
        \draw[-{Stealth}, thick, PoliBlue!80] (0,0) -- (3.5,0) node[below right, text=PoliBlue] {$t$};
        \draw[thick, orange!70!yellow, domain=0:3, samples=100] plot (\x, {cos(\x*360/2)});
        \draw[|-|, thin, gray] (-0.2, 0) -- (-0.2, 1) node[midway, left, text=PoliBlue] {$A$};
        \draw[|-|, thin, gray] (0, 1.2) -- (2, 1.2) node[midway, above, text=orange] {$T$};
        
        % --- FRECCIA DI CONVERSIONE ---
        \draw[-{Stealth[scale=1.5]}, line width=4pt, PoliBlue!70] (4,0) -- (5,0);
        
        % --- DOMINIO DELLA FREQUENZA (Destra) ---
        \begin{scope}[shift={(5.5, -1.5)}]
            \draw[-{Stealth}, thick, red] (0,0) -- (0,3) node[right] {$P$};
            \draw[-{Stealth}, thick, red] (0,0) -- (4,0) node[above right] {$\nu$};
            
            % Spettro di potenza (Singolo picco)
            \draw[thick, ForestGreen] (2, 0) -- (2, 1.5);
            \draw[thick, ForestGreen] (1.95, 0) -- (2, 1.5) -- (2.05, 0); % Effetto "delta"
        \end{scope}
    \end{tikzpicture}
\end{figure}

\begin{itemize}[label=\textcolor{PoliBlue}{\textbullet}]
    \item Un segnale sinusoidale periodico che avrà periodo più breve $T_2$ rispetto a $T_1$ avrà frequenza $\nu_2 > \nu_1$.
\end{itemize}

\begin{figure}[H]
    \centering
    \begin{tikzpicture}[xscale=1.2, yscale=1.2]
        % --- DOMINIO DEL TEMPO (Sinistra) ---
        \draw[-{Stealth}, thick, PoliBlue] (0,-2) -- (0,2) node[left] {s};
        \draw[-{Stealth}, thick, PoliBlue] (0,0) -- (4.5,0) node[below right] {t};
        
        % Curva T1 (Arancione, periodo lungo)
        \draw[thick, orange, domain=0:4, samples=100] plot (\x, {1.2*cos(\x*360/3)});
        
        % Curva T2 (Giallo, periodo corto)
        \draw[thick, orange!60!yellow, domain=0:4, samples=150] plot (\x, {1.2*cos(\x*360/1)});
        
        % Quote
        \draw[|-|, thin, gray] (-0.3, 0) -- (-0.3, 1.2) node[midway, left, text=PoliBlue] {A};
        \draw[|-|, thin, orange] (0, 1.5) -- (3, 1.5) node[midway, above] {$T_1$};
        \draw[|-|, thin, orange!60!yellow] (1, -1.5) -- (2, -1.5) node[midway, below] {$T_2$};
        
        % --- FRECCIA DI CONVERSIONE ---
        \draw[-{Stealth[scale=1.5]}, line width=4pt, PoliBlue!70] (5,0) -- (5.8,0);
        
        % --- DOMINIO DELLA FREQUENZA (Destra) ---
        \begin{scope}[shift={(6.5, -2)}]
            \draw[-{Stealth}, thick, red] (0,0) -- (0,4) node[left] {P};
            \draw[-{Stealth}, thick, red] (0,0) -- (5,0) node[above right] {$\nu$};
            
            % Picco T1 (Frequenza v1 minore -> a sinistra)
            \draw[thick, orange] (1.5, 0) -- (1.5, 2);
            \draw[thick, orange] (1.45, 0) -- (1.5, 2) -- (1.55, 0);
            \node[below, text=orange] at (1.5, -0.2) {$\nu_1$};
            
            % Picco T2 (Frequenza v2 maggiore -> a destra)
            \draw[thick, orange!60!yellow] (3.5, 0) -- (3.5, 2);
            \draw[thick, orange!60!yellow] (3.45, 0) -- (3.5, 2) -- (3.55, 0);
            \node[below, text=orange!60!yellow] at (3.5, -0.2) {$\nu_2$};
        \end{scope}
    \end{tikzpicture}
\end{figure}

\begin{itemize}[label=\textcolor{PoliBlue}{\textbullet}]
    \item Un segnale sinusoidale periodico che avrà ampiezza maggiore $A_2$ rispetto a $A_1$ avrà potenza $P_2 > P_1$.
\end{itemize}

\begin{figure}[H]
    \centering
    \begin{tikzpicture}[xscale=1.2, yscale=1.2]
        % --- DOMINIO DEL TEMPO (Sinistra) ---
        \draw[-{Stealth}, thick, PoliBlue] (0,-2.2) -- (0,2.2) node[left] {s};
        \draw[-{Stealth}, thick, PoliBlue] (0,0) -- (4.5,0) node[below right] {t};
        
        % Curva 1 (Arancione, Ampiezza minore A1, stesso periodo T)
        \draw[thick, orange, domain=0:4, samples=100] plot (\x, {1.2*cos(\x*360/3.2)});
        
        % Curva 2 (Grigia, Ampiezza maggiore A2, stesso periodo T)
        \draw[thick, gray!70, domain=0:4, samples=100] plot (\x, {1.8*cos(\x*360/3.2)});
        
        % Quote Ampiezze
        \draw[|-|, thin, orange] (-0.3, 0) -- (-0.3, 1.2) node[midway, left] {$A_1$};
        \draw[|-|, thin, gray!70] (-0.7, 0) -- (-0.7, 1.8) node[midway, left] {$A_2$};
        
        % Quota Periodo (Condivisa)
        \draw[|-|, thin, PoliBlue] (0.8, 2) -- (4.0, 2) node[midway, above] {T};
        
        % --- FRECCIA DI CONVERSIONE ---
        \draw[-{Stealth[scale=1.5]}, line width=4pt, PoliBlue!70] (4.8,0) -- (5.6,0);
        
        % --- DOMINIO DELLA FREQUENZA (Destra) ---
        \begin{scope}[shift={(6.2, -2.2)}]
            \draw[-{Stealth}, thick, red] (0,0) -- (0,4.4) node[left] {P};
            \draw[-{Stealth}, thick, red] (0,0) -- (5,0) node[above right] {$\nu$};
            
            % Etichette Assi P
            \node[left, text=orange] at (0, 1.7) {$P_1$};
            \node[left, text=gray!80]  at (0, 3) {$P_2$};
            
            % Picchi della trasformata. 
            % Essendo la stessa frequenza (stesso T) i picchi cadono nello stesso punto.
            % Li disegniamo leggermente sfalsati solo per renderli entrambi visibili, come nella slide.
            
            % Picco A1 (Arancione, Potenza minore)
            \draw[thick, orange] (2.35, 0) -- (2.35, 1.7);
            \draw[thick, orange] (2.3, 0) -- (2.35, 1.7) -- (2.4, 0);
            
            % Picco A2 (Grigio, Potenza maggiore)
            \draw[thick, gray!80] (2.45, 0) -- (2.45, 3);
            \draw[thick, gray!80] (2.4, 0) -- (2.45, 3) -- (2.5, 0);
            
            % Etichetta asse X (Frequenze uguali)
            \node[below] at (2.4, -0.2) {\textcolor{orange}{$\nu_1$} = \textcolor{gray!80}{$\nu_2$}};
        \end{scope}
    \end{tikzpicture}
\end{figure}

\begin{itemize}[label=\textcolor{PoliBlue}{\textbullet}]
    \item E naturalmente valgono anche tutte le varie combinazioni...
\end{itemize}

\begin{figure}[H]
    \centering
    \begin{tikzpicture}[xscale=1.2, yscale=1.2]
        % --- DOMINIO DEL TEMPO (Sinistra) ---
        \draw[-{Stealth}, thick, PoliBlue] (0,-2) -- (0,2) node[left] {s};
        \draw[-{Stealth}, thick, PoliBlue] (0,0) -- (5,0) node[below right] {t};
        
        % Curva 1 (Arancione, Ampiezza alta, Periodo corto T1)
        \draw[thick, orange, domain=0:4.5, samples=100] plot (\x, {1.2*cos(\x*360/2.5)});
        
        % Curva * (Nera, Ampiezza bassa, Periodo lungo T*)
        \draw[thick, black, domain=0:4.5, samples=100] plot (\x, {0.6*cos(\x*360/4.5)});
        
        % Quote Ampiezze
        \draw[|-|, thin, orange] (-0.5, 0) -- (-0.5, 1.2) node[midway, left] {$A_1$};
        \draw[|-|, thin, black] (-0.2, 0) -- (-0.2, 0.6) node[midway, left] {$A_*$};
        
        % Quote Periodi
        \draw[|-|, thin, orange] (1.25, 1.5) -- (3.75, 1.5) node[midway, above] {$T_1$};
        \draw[|-|, thin, black] (0, -1.8) -- (4.5, -1.8) node[midway, above] {$T_*$};
        
        % --- FRECCIA DI CONVERSIONE ---
        \draw[-{Stealth[scale=1.5]}, line width=4pt, PoliBlue!70] (5.3,0) -- (6.1,0);
        
        % --- DOMINIO DELLA FREQUENZA (Destra) ---
        \begin{scope}[shift={(6.8, -2)}]
            \draw[-{Stealth}, thick, red] (0,0) -- (0,4) node[left] {P};
            \draw[-{Stealth}, thick, red] (0,0) -- (5,0) node[above right] {$\nu$};
            
            % Etichette Assi P
            \node[left, text=orange] at (0, 2) {$P_1$};
            \node[left, text=black]  at (0, 1) {$P_*$};
            
            % Picco * (Periodo T* lungo -> Frequenza bassa -> a sinistra. Ampiezza bassa -> Potenza bassa)
            \draw[thick, black] (1, 0) -- (1, 1);
            \draw[thick, black] (0.95, 0) -- (1, 1) -- (1.05, 0);
            \node[below, text=black] at (1, -0.2) {$\nu_*$};
            
            % Picco 1 (Periodo T1 corto -> Frequenza alta -> a destra. Ampiezza alta -> Potenza alta)
            \draw[thick, orange] (2.5, 0) -- (2.5, 2);
            \draw[thick, orange] (2.45, 0) -- (2.5, 2) -- (2.55, 0);
            \node[below, text=orange] at (2.5, -0.2) {$\nu_1$};
        \end{scope}
    \end{tikzpicture}
\end{figure}

\subsection*{Composizione dei Segnali e Teoria di Fourier}

Finora stiamo trattando segnali composti da una sola sinusoide: è bene ricordare che per la verità i segnali reali sono costituiti da varie componenti che coprono una varietà di frequenze. La teoria di Fourier ci aiuta a individuare le varie componenti.

\begin{figure}[H]
    \centering
    \begin{tikzpicture}[xscale=1.5, yscale=2, font=\sffamily\small]
        
        % Griglia e assi
        \draw[gray!30, thin] (0,-1.5) grid (3*pi,1.5);
        \draw[-{Stealth}, thin, gray] (0,0) -- (3*pi+0.5,0);
        \draw[-{Stealth}, thin, gray] (0,-1.5) -- (0,1.5);
        
        % Labels asse Y
        \foreach \y in {-1.5, -1, -0.5, 0, 0.5, 1, 1.5}
            \node[left, text=gray, font=\tiny] at (0,\y) {\y};
            
        % Fascia Armoniche (Rettangolo viola in background)
        \fill[blue!30!purple, opacity=0.3] (0, -0.45) rectangle (3*pi, 0.45);
        
        % Sintesi onda quadra ideale (Teal scuro)
        \draw[thick, teal, opacity=0.7] (0,1) -- (pi,1) -- (pi,-1) -- (2*pi,-1) -- (2*pi,1) -- (3*pi,1);
        
        % Fondamentale (Blu)
        \draw[thick, blue, domain=0:3*pi, samples=150] plot (\x, {sin(\x r)});
        
        % Armoniche dispari (Alcune tracciate per simulare la fascia)
        \draw[purple, thin, domain=0:3*pi, samples=200] plot (\x, {1/3*sin(3*\x r)});
        \draw[ForestGreen, thin, domain=0:3*pi, samples=200] plot (\x, {1/5*sin(5*\x r)});
        \draw[orange, thin, domain=0:3*pi, samples=200] plot (\x, {1/7*sin(7*\x r)});
        \draw[cyan, thin, domain=0:3*pi, samples=200] plot (\x, {1/9*sin(9*\x r)});
        
        % Somma (Arossata, approssimazione con le prime armoniche per simulare il grafico)
        \draw[thick, red!60, domain=0:3*pi, samples=300] plot (\x, {sin(\x r) + 1/3*sin(3*\x r) + 1/5*sin(5*\x r) + 1/7*sin(7*\x r) + 1/9*sin(9*\x r) + 1/11*sin(11*\x r)});
        
        % --- SEZIONE MODIFICATA: Frecce e Testo ---
        
        % Fondamentale
        \node[text=blue, font=\bfseries] (fond) at (1.8, 1.25) {Fondamentale};
        \draw[-{Stealth[scale=1.2]}, ultra thick, blue] (fond.south east) -- (2.3, 0.75);
        
        % Somma
        \node[text=blue, font=\bfseries] (somma) at (3.3, 0.85) {Somma};
        \draw[-{Stealth[scale=1.2]}, ultra thick, red] (somma.south west) -- (3.0, 0.45);
        
        % Armoniche dispari
        \node[text=blue, font=\bfseries, align=left] (arm) at (6.5, 1.0) {Armoniche\\dispari};
        \draw[-{Stealth[scale=1.2]}, ultra thick, gray!90!black] (arm.south west) -- (5.6, 0.15);
        
        % ------------------------------------------
        
        % Titolo inferiore
        \node[below, text=red, font=\bfseries\large] at (1.5*pi, -1.1) {Sintesi onda quadra (12 armoniche)};
        
    \end{tikzpicture}
\end{figure}


\section{Il Rumore}

\subsection*{Introduzione al Rumore}

Tra i vari motivi per cui quanto trattato finora corrisponde a un segnale ideale, c'è il fatto che non abbiamo considerato il rumore.

\begin{itemize}[label=\textcolor{PoliBlue}{\textbullet}]
    \item Nella realtà abbiamo \textbf{sempre} almeno un \textbf{rumore di fondo}, anche se di piccola entità, cioè ampiezza, ovvero potenza.
    \begin{itemize}[label=$\circ$]
        \item Per esempio, quando una persona parla, esiste sempre un rumore di fondo al di sopra del quale si cerca di alzare il volume della voce.
        \item Inoltre, il rumore è introdotto \textbf{non solo dall'ambiente}, ma dagli apparecchi stessi che rilevano il segnale.
    \end{itemize}
\end{itemize}

Il rumore consiste in variazioni indesiderate del segnale originate nel segnale stesso o dal sistema di comunicazione:
\begin{itemize}[label=\textcolor{OpticRed}{\textbullet}]
    \item \textbf{non può essere evitato;}
    \item \textbf{però può essere analizzato e mitigato.}
\end{itemize}

\subsection*{Rappresentazione del Rumore}

Il rumore può essere rappresentato come una composizione di altre sinusoidi di varie lunghezza d'onda che si sovrappongono indistinguibilmente tra loro.

\begin{figure}[H]
    \centering
    \begin{tikzpicture}[xscale=1.2, yscale=1.2]
        % --- DOMINIO DEL TEMPO (Sinistra) ---
        \draw[-{Stealth}, thick, PoliBlue] (0,-1.5) -- (0,1.8) node[left] {A};
        \draw[-{Stealth}, thick, PoliBlue] (0,0) -- (4.5,0) node[below right] {t};
        
        % Segnale principale
        \draw[thick, orange, domain=0:4, samples=100] plot (\x, {1.2*cos(\x*360/3)});
        
        % Rumore (sovrapposizione di sinusoidi piccole)
        \draw[thick, black, domain=0:4, samples=100] plot (\x, {0.15*cos(\x*360/1.8) + 0.1});
        \draw[thick, ForestGreen, domain=0:4, samples=150] plot (\x, {0.2*sin(\x*360/2.1) - 0.05});
        \draw[thick, orange!40, domain=0:4, samples=150] plot (\x, {0.15*cos(\x*360/1.1) - 0.15});
        
        % --- DOMINIO DELLA FREQUENZA (Destra) ---
        \begin{scope}[shift={(6.5, -1.5)}]
            \draw[-{Stealth}, thick, red] (0,0) -- (0,3) node[left] {P};
            \draw[-{Stealth}, thick, red] (0,0) -- (5,0) node[above right] {$\nu$};
            
            % Picco Segnale Principale
            \draw[thick, orange] (2, 0) -- (2, 2.5);
            \draw[thick, orange] (1.95, 0) -- (2, 2.5) -- (2.05, 0);
            
            % Picchi Rumore
            \foreach \x/\h in {0.5/0.2, 0.8/0.25, 1.2/0.2, 2.8/0.25, 3.2/0.2, 3.8/0.25} {
                \draw[thick, black] (\x, 0) -- (\x, \h);
                \draw[thick, black] (\x-0.03, 0) -- (\x, \h) -- (\x+0.03, 0);
            }
        \end{scope}
    \end{tikzpicture}
\end{figure}

Il rumore, all'atto pratico, si somma al segnale alterandolo visivamente.

\begin{figure}[H]
    \centering
    \begin{tikzpicture}[xscale=1.1, yscale=1.1]

        % Seed per avere un rumore casuale ma riproducibile
        \pgfmathsetseed{12345}

        % --- DOMINIO DEL TEMPO (Sinistra) ---
        \draw[-{Stealth}, thick, PoliBlue] (0,-1.5) -- (0,1.8) node[left] {A};
        \draw[-{Stealth}, thick, PoliBlue] (0,0) -- (4,0) node[below right] {t};

        % Segnale + rumore
        % Il segnale sinusoidale viene perturbato da piccole variazioni casuali
        \draw[
            thick,
            orange,
            domain=0:3.8,
            samples=180
        ]
        plot (\x, {
            1.2*cos(\x*360/3)
            + 0.12*rand
        });

        % --- DOMINIO DELLA FREQUENZA (Destra) ---
        \begin{scope}[shift={(6,-1.5)}]

            \draw[-{Stealth}, thick, red]
                (0,0) -- (0,3) node[left] {P};

            \draw[-{Stealth}, thick, red]
                (0,0) -- (4.5,0) node[above right] {$\nu$};

            % Fondo di rumore casuale
            \draw[
                thick,
                black!80
            ]
            plot[
                domain=0.05:4,
                samples=80
            ]
            (\x,{0.10 + 0.12*rand});

            % Picco del segnale
            \draw[thick, orange]
                (2,0) -- (2,2.5);

            \draw[thick, orange]
                (1.95,0) -- (2,2.5) -- (2.05,0);

        \end{scope}

    \end{tikzpicture}
\end{figure}
Il problema critico sorge quando il segnale è molto basso o il rumore è molto alto: in queste condizioni \textbf{non si distingue più il segnale dal rumore}.

\begin{figure}[H]
    \centering
    \begin{tikzpicture}[xscale=1.1, yscale=1.1]

        % Seed diverso per il secondo grafico
        \pgfmathsetseed{67890}

        % --- DOMINIO DEL TEMPO (Sinistra) ---
        \draw[-{Stealth}, thick, PoliBlue]
            (0,-1) -- (0,1.5) node[left] {A};

        \draw[-{Stealth}, thick, PoliBlue]
            (0,0) -- (4.5,0) node[below right] {t};

        % Segnale debole, mostrato separatamente
        \draw[
            thick,
            orange!70,
            domain=0:4,
            samples=150
        ]
        plot (\x,{0.30*cos(\x*360/3)});

        % Segnale osservato: segnale + rumore
        % Il rumore ora è molto più grande
        \draw[
            thick,
            black,
            domain=0:4,
            samples=180
        ]
        plot (\x,{
            0.30*cos(\x*360/3)
            + 0.32*rand
        });

        % --- DOMINIO DELLA FREQUENZA (Destra) ---
        \begin{scope}[shift={(6,-1)}]

            \draw[-{Stealth}, thick, red]
                (0,0) -- (0,2.5) node[left] {P};

            \draw[-{Stealth}, thick, red]
                (0,0) -- (4.5,0) node[above right] {$\nu$};

            % Rumore elevato
            \draw[
                thick,
                black!80
            ]
            plot[
                domain=0.05:4,
                samples=80
            ]
            (\x,{0.12 + 0.40*rand});

            % Picco del segnale debole:
            % ora è appena distinguibile dal rumore
            \draw[thick, orange]
                (2,0) -- (2,0.65);

            \draw[thick, orange]
                (1.95,0) -- (2,0.65) -- (2.05,0);

        \end{scope}

    \end{tikzpicture}
\end{figure}

\subsection*{Rapporto S/R o SNR}

Dunque, ciò che si desidera avere lungo tutto il percorso di un segnale (che sia un segnale naturale che stiamo captando o che sia un segnale artificiale trasmesso e ricevuto) è un migliore possibile rapporto segnale-rumore, indicato con \textbf{S/R} o \textbf{SNR} (\textit{Signal to Noise Ratio}).

\begin{itemize}[label=\textcolor{PoliBlue}{\textbullet}]
    \item Questo termine indica il \textbf{rapporto} tra il livello di \textbf{potenza} del \textbf{segnale} ricevuto $S$ e il livello di \textbf{potenza del rumore} $N$.
    \item Essendo un rapporto di potenze, possiamo esprimerlo elegantemente come differenza in decibel (dB).
\end{itemize}

\begin{equation}
    \text{SNR} = \frac{S}{N} = \frac{P_S}{P_N}
\end{equation}

\begin{figure}[H]
    \centering
    \begin{tikzpicture}[xscale=1.5, yscale=1.5]

        % Seed per avere un rumore casuale ma riproducibile
        \pgfmathsetseed{24680}

        % =========================================================
        % ASSI
        % =========================================================

        \draw[-{Stealth}, ultra thick, red]
            (0,0) -- (0,3)
            node[left] {$P$};

        \draw[-{Stealth}, ultra thick, red]
            (0,0) -- (5,0)
            node[below right] {$\nu$};


        % =========================================================
        % RUMORE - NOISE FLOOR
        % =========================================================

        \draw[
            thick,
            black!80
        ]
        plot[
            domain=0.05:4.7,
            samples=160
        ]
        (\x,{
            0.12
            + 0.06*rand
            + 0.04*sin(\x*900)
        });


        % =========================================================
        % PICCO DEL SEGNALE
        % =========================================================

        % Parte centrale del picco
        \draw[
            thick,
            orange
        ]
        plot[
            domain=2.43:2.57,
            samples=40
        ]
        (\x,{
            1.45*exp(-((\x-2.5)/0.025)^2)
        });


        % =========================================================
        % LIVELLI Ps E Pn
        % =========================================================

        % Linea verticale di riferimento
        \draw[
            thick,
            PoliBlue,
            <-> 
        ]
        (-0.75,0.12) -- (-0.75,1.45);

        % Piccola separazione visiva tra P_N e P_S
        \draw[
            thick,
            PoliBlue
        ]
        (-0.82,0.12) -- (-0.68,0.12);

        \draw[
            thick,
            PoliBlue
        ]
        (-0.82,1.45) -- (-0.68,1.45);


        % =========================================================
        % ETICHETTE
        % =========================================================

        \node[
            text=orange,
            font=\bfseries
        ]
        at (-0.82,1.45)
        [left]
        {$P_S = S$};

        \node[
            text=black,
            font=\bfseries
        ]
        at (-0.82,0.12)
        [left]
        {$P_N = N$};

    \end{tikzpicture}
\end{figure}

\subsection*{Sorgenti e Tipi di Rumore}

La potenza del rumore è il risultato di processi casuali come il flusso di cariche o di lacune in un dispositivo a stato solido, la propagazione attraverso la ionosfera o altri gas ionizzati o le vibrazioni termiche di qualsiasi componente a una temperatura superiore allo zero assoluto. Quindi le sorgenti del rumore possono essere:

\begin{itemize}[label=\textcolor{PoliBlue}{\textbullet}]
    \item emesse da \textbf{sorgenti esterne} e che sono raccolte dal ricevitore (cioè dall'area dell'antenna: atmosfera, spazio, \dots),
    \item generate all'\textbf{interno} dai componenti stessi del sistema ricevente.
\end{itemize}

Tali disturbi interni possono essere dovuti a diversi processi, dando origine a vari tipi di rumore:

\begin{itemize}[label=\textcolor{PoliBlue}{\textbullet}]
    \item Il \textbf{rumore termico}, noto anche come rumore di \textbf{Johnson} o di \textbf{Nyquist}, è il tipo più elementare di rumore, causato da movimenti/\textbf{vibrazioni} casuali di cariche o portatori di carica nei dispositivi e nei materiali:
    \begin{itemize}[label=$\circ$]
        \item è prodotto in tutti i dispositivi conduttori: perciò ogni sistema di comunicazione aggiunge rumore;
    \end{itemize}
    
    \item Lo \textbf{shot noise} è dovuto alle fluttuazioni casuali dei portatori di carica;
    
    \item Il \textbf{rumore di sfarfallamento} (\textit{flicker noise} o \textit{pink noise}) si verifica nei componenti in cui la potenza del rumore varia inversamente con la frequenza, per cui viene spesso chiamato \textbf{rumore 1/f} (quindi interessa le basse frequenze, sotto \qty{1}{\kilo\hertz});
    
    \item Il \textbf{rumore del plasma} è causato dal movimento casuale delle cariche in un \textbf{gas ionizzato};
    
    \item Il \textbf{rumore quantistico} deriva dalla natura quantizzata dei portatori di carica e dei fotoni; spesso è insignificante rispetto ad altre fonti di rumore.
\end{itemize}


\chapter{Sistemi di Telecomunicazioni}

Un sistema di comunicazioni muove le informazioni da una sorgente a una destinazione attraverso un canale, cercando di non alterarne il contenuto.

In particolare, una radio opera su segnali elettromagnetici, agendo sulla loro \textbf{frequenza} e la \textbf{potenza}, per generare un segnale/frequenza portante o semplicemente portante, che possa essere propagata attraverso il canale desiderato.

\vspace{0.5cm}

\begin{figure}[H]
    \centering
    \resizebox{\textwidth}{!}{% Ridimensiona per far stare tutti i blocchi nella larghezza della pagina
    \begin{tikzpicture}[
        block/.style={rectangle, draw=PoliBlue, thick, minimum height=1.2cm, align=center, text width=2.2cm, fill=white, font=\sffamily\small},
        arrow/.style={-{Stealth}, thick, PoliBlue}
    ]
        % Nodi (Blocchi)
        \node[block] (sorg) {sorgente};
        \node[block, right=0.6cm of sorg] (trasd1) {trasduttore};
        \node[block, right=0.6cm of trasd1] (tx) {trasmettitore};
        \node[block, right=0.6cm of tx] (canale) {canale};
        \node[block, right=0.6cm of canale] (rx) {ricevitore};
        \node[block, right=0.6cm of rx] (trasd2) {detector/\\trasduttore};
        \node[block, right=0.6cm of trasd2] (dest) {destinazione};

        % Frecce di collegamento
        \draw[arrow] (sorg) -- (trasd1);
        \draw[arrow] (trasd1) -- (tx);
        \draw[arrow] (tx) -- (canale);
        \draw[arrow] (canale) -- (rx);
        \draw[arrow] (rx) -- (trasd2);
        \draw[arrow] (trasd2) -- (dest);

        % Indicatore tratteggiato per il "sistema di telecomunicazioni"
        \coordinate (mid_sys) at ($(tx.south)!0.5!(rx.south)$);
        \coordinate (label_pos) at ($(mid_sys) - (0, 1.8cm)$);

        \draw[dashed, thick, gray!80] (tx.south) -- (label_pos) -- (rx.south);
        \node[below=0.1cm of label_pos, font=\rmfamily\large, text=blue!90!black] {sistema di telecomunicazioni};

    \end{tikzpicture}
    }
\end{figure}

\vspace{0.5cm}
\begin{itemize}[label=\textcolor{gray}{\textbullet}]
    \item \textcolor{gray}{(qui ci occuperemo solo dei segmenti della radio; successivamente del canale)}
\end{itemize}

\section{Banda Base e Banda Passante}
\subsection*{Banda Base}
Ragioniamo d'ora in poi in termini di frequenze. In molti casi i segnali prodotti hanno frequenze basse:
\begin{itemize}[label=\textcolor{PoliBlue}{\textbullet}]
    \item pensiamo alla voce umana che viene emessa e udita tra i \qty{20}{\hertz} e i \qty{20}{\kilo\hertz}.
\end{itemize}

\begin{figure}[H]
    \centering
    \begin{tikzpicture}[xscale=1.5, yscale=1.5, font=\sffamily\small]
        % Assi
        \draw[-{Stealth}, thin] (0,0) -- (5,0) node[below] {\scriptsize Frequenza};
        \draw[-{Stealth}, thin] (0.5,-0.2) -- (0.5,2);
        
        % Etichetta Asse Y ruotata
        \node[rotate=90, anchor=south] at (0.2, 1) {\scriptsize Potenza di uscita};
        
        % Tacche e Testi asse X
        \draw (1.5, 0.1) -- (1.5, -0.1) node[below] {\scriptsize \qty{20}{\hertz}};
        \node[above] at (1.5, 0.1) {\scriptsize $f_L$};
        
        \draw (3.5, 0.1) -- (3.5, -0.1) node[below] {\scriptsize \qty{20}{\kilo\hertz}};
        \node[above] at (3.5, 0.1) {\scriptsize $f_H$};
        
        % Tacche e Testi asse Y
        \draw (0.4, 1.5) -- (0.6, 1.5) node[left, xshift=-0.2cm] {\scriptsize $P_{\text{mid}}$};
        
        % Curva a trapezio (Risposta dell'amplificatore)
        \draw[thick] (0.8, 0.3) -- (1.5, 1.5) -- (3.5, 1.5) -- (4.2, 0.3);
    \end{tikzpicture}
    \caption*{Potenza di uscita in funzione della frequenza del segnale di ingresso di un amplificatore audio.}
\end{figure}

Anche altri segnali vengono prodotti nella cosiddetta \textbf{banda base}, ovvero alla loro frequenza naturale relativamente molto vicina alla frequenza \textbf{\qty{0}{\hertz}}:
\begin{itemize}[label=\textcolor{PoliBlue}{\textbullet}]
    \item ovvero dal valore di \qty{0}{\hertz} al valore di $\nu_S$ (frequenza limite stabilita dalla forma d'onda del segnale):
    \begin{itemize}[label=$\circ$]
        \item un segnale video analogico di un broadcast televisivo va da corrente continua (o DC) a \qty{4.2}{\mega\hertz}.
    \end{itemize}
\end{itemize}

\subsection*{Definizioni di Banda}
Consideriamo un segnale $s(t)$ con estremi di banda $(f_1; f_2)$: dunque la \textbf{larghezza di banda} è $B = f_2 - f_1$. \\
Inoltre, denotiamo con $f_0$ la \textbf{frequenza di centro banda} di $s(t)$, cioè:
\begin{equation}
    f_0 = \frac{f_1 + f_2}{2}
\end{equation}

La \textbf{larghezza di banda relativa} è definita come $B / f_0$.

\begin{nota}
    Allora la \textbf{banda base} si ha quando $B \approx 2f_0$, oppure equivalentemente quando la larghezza di banda relativa è prossima al valore 2 (cioè, $B / f_0 \approx 2$).
\end{nota}

Si parla invece di \textbf{banda passante} o \textbf{banda traslata} quando la stessa porzione di banda di frequenze è spostata a \textbf{frequenze maggiori}, e dunque in tal caso avremo:
\begin{equation}
    \frac{B}{f_0} \ll 1
\end{equation}

\vspace{0.5cm}

\begin{figure}[H]
    \centering
    \begin{tikzpicture}[
        xscale=1.1,
        yscale=1.2,
        font=\sffamily\small
    ]

    % Seed per rendere il rumore riproducibile
    \pgfmathsetseed{31415}

    % ============================================================
    % GRAFICO 1 - BANDA BASE
    % ============================================================

    \begin{scope}[shift={(0,0)}]

        % --------------------------------------------------------
        % Assi
        % --------------------------------------------------------

        \draw[-{Stealth}, ultra thick, red]
            (0,0) -- (0,2.5)
            node[left] {$P$};

        \draw[-{Stealth}, ultra thick, red]
            (0,0) -- (4.2,0)
            node[below right] {$\nu$};


        % --------------------------------------------------------
        % Rumore
        % --------------------------------------------------------

        \draw[
            thick,
            black!80
        ]
        plot[
            domain=0.05:4.0,
            samples=180
        ]
        (\x,{
            0.12
            + 0.07*rand
            + 0.025*sin(900*\x)
        });


        % --------------------------------------------------------
        % Banda del segnale
        % --------------------------------------------------------

        \draw[
            thick,
            orange
        ]
        (0.20,0.12)
        -- (0.40,1.45)
        -- (1.55,1.45)
        -- (1.75,0.12);


        % --------------------------------------------------------
        % Frequenze f1, f0, f2
        % --------------------------------------------------------

        % f1
        \draw[
            thin,
            orange,
            dashed
        ]
        (0.40,0) -- (0.40,-0.45);

        \node[
            text=orange
        ]
        at (0.40,-0.52)
        {$f_1$};


        % f0
        \draw[
            thin,
            orange,
            dashed
        ]
        (0.975,0) -- (0.975,-0.45);

        \node[
            text=orange
        ]
        at (0.975,-0.52)
        {$f_0$};


        % f2
        \draw[
            thin,
            orange,
            dashed
        ]
        (1.55,0) -- (1.55,-0.45);

        \node[
            text=orange
        ]
        at (1.55,-0.52)
        {$f_2$};


        % --------------------------------------------------------
        % Larghezza di banda B
        % --------------------------------------------------------

        \draw[
            <->,
            thick,
            orange
        ]
        (0.40,-0.85)
        -- (1.55,-0.85);

        \node[
            text=orange
        ]
        at (0.975,-1.05)
        {$B$};

    \end{scope}


    % ============================================================
    % GRAFICO 2 - BANDA PASSANTE
    % ============================================================

    \begin{scope}[shift={(5.8,0)}]

        % --------------------------------------------------------
        % Assi
        % --------------------------------------------------------

        \draw[-{Stealth}, ultra thick, red]
            (0,0) -- (0,2.5)
            node[left] {$P$};

        \draw[-{Stealth}, ultra thick, red]
            (0,0) -- (4.5,0)
            node[below right] {$\nu$};


        % --------------------------------------------------------
        % Rumore
        % --------------------------------------------------------

        \draw[
            thick,
            black!80
        ]
        plot[
            domain=0.05:4.3,
            samples=190
        ]
        (\x,{
            0.12
            + 0.07*rand
            + 0.025*sin(850*\x)
        });


        % --------------------------------------------------------
        % Banda del segnale traslata
        % --------------------------------------------------------

        \draw[
            thick,
            orange
        ]
        (1.80,0.12)
        -- (2.00,1.45)
        -- (3.15,1.45)
        -- (3.35,0.12);


        % --------------------------------------------------------
        % Frequenze f1, f0, f2
        % --------------------------------------------------------

        % f1
        \draw[
            thin,
            orange,
            dashed
        ]
        (2.00,0) -- (2.00,-0.45);

        \node[
            text=orange
        ]
        at (2.00,-0.52)
        {$f_1$};


        % f0
        \draw[
            thin,
            orange,
            dashed
        ]
        (2.575,0) -- (2.575,-0.45);

        \node[
            text=orange
        ]
        at (2.575,-0.52)
        {$f_0$};


        % f2
        \draw[
            thin,
            orange,
            dashed
        ]
        (3.15,0) -- (3.15,-0.45);

        \node[
            text=orange
        ]
        at (3.15,-0.52)
        {$f_2$};


        % --------------------------------------------------------
        % Larghezza di banda B
        % --------------------------------------------------------

        \draw[
            <->,
            thick,
            orange
        ]
        (2.00,-0.85)
        -- (3.15,-0.85);

        \node[
            text=orange
        ]
        at (2.575,-1.05)
        {$B$};

    \end{scope}

    \end{tikzpicture}
\end{figure}

\subsection*{Banda Passante}

Se pensiamo alle usuali telecomunicazioni, queste non occupano / si trovano a bassa frequenza direttamente in banda base, quanto piuttosto ad alta frequenza, ovvero in \textbf{banda passante}:
\begin{itemize}[label=\textcolor{PoliBlue}{\textbullet}]
    \item per esempio è noto che le trasmissioni radiofoniche avvengono in canali di frequenza attorno i \qty{100}{\mega\hertz};
    \item se pensiamo alle frequenze in uso per comunicare con lo spazio, queste sono oltre il \qty{1}{\giga\hertz}.
\end{itemize}

\begin{table}[H]
    \centering
    \small
    \renewcommand{\arraystretch}{1.1}
    \begin{tabular}{lc}
        \toprule
        \textbf{Categorie} & \textbf{Gamma di frequenza} \\
        \midrule
        % Evidenziazione in verde per la banda base (come il rettangolo nella slide)
        \rowcolor{ForestGreen!10} \textcolor{ForestGreen}{\textbf{Suoni udibili}} & \textcolor{ForestGreen}{\qty{20}{\hertz} $\div$ \qty{20}{\kilo\hertz}} \\
        \rowcolor{ForestGreen!10} \textcolor{ForestGreen}{\textbf{Segnale video TV in banda di base}} & \textcolor{ForestGreen}{\qty{0}{\hertz} $\div$ \qty{4.5}{\mega\hertz}} \\
        Diffusione radio AM & \qty{540}{\kilo\hertz} $\div$ \qty{1600}{\kilo\hertz} \\
        Comunicazioni radio ad alta frequenza & \qty{1.6}{\mega\hertz} $\div$ \qty{54}{\mega\hertz} \\
        Televisione VHF (canali 2-6) & \qty{54}{\mega\hertz} $\div$ \qty{88}{\mega\hertz} \\
        \midrule
        Diffusione radio FM & \qty{88}{\mega\hertz} $\div$ \qty{108}{\mega\hertz} \\
        Comunicazione radio VHF & \qty{108}{\mega\hertz} $\div$ \qty{174}{\mega\hertz} \\
        Televisione VHF (canali 7-13) & \qty{174}{\mega\hertz} $\div$ \qty{216}{\mega\hertz} \\
        \midrule
        Comunicazioni marittime e governative & \qty{216}{\mega\hertz} $\div$ \qty{450}{\mega\hertz} \\
        Comunicazioni d'affari & \qty{450}{\mega\hertz} $\div$ \qty{470}{\mega\hertz} \\
        Televisione UHF (canali 14-69) & \qty{470}{\mega\hertz} $\div$ \qty{806}{\mega\hertz} \\
        Allocazione di banda per comunicazioni fisse e mobili & \qty{806}{\mega\hertz} $\div$ \qty{902}{\mega\hertz} \\
        Allocazione di banda per cellulari analogici e digitali & \qty{928}{\mega\hertz} $\div$ \qty{960}{\mega\hertz} \\
        Telefonia, comunicazioni personali e altro & \qty{1710}{\mega\hertz} $\div$ \qty{1990}{\mega\hertz} \\
        Dispositivi wireless & \qty{2310}{\mega\hertz} $\div$ \qty{2690}{\mega\hertz} \\
        Televisione satellitare & \qty{3.7}{\giga\hertz} $\div$ \qty{4.2}{\giga\hertz} \\
        Dispositivi wireless & \qty{5.0}{\giga\hertz} $\div$ \qty{5.5}{\giga\hertz} \\
        \bottomrule
    \end{tabular}
    \caption*{Frequenze associate a segnali di impiego comune.}
\end{table}



\section{Alta Frequenza e Teorema di Shannon}

Le \textbf{motivazioni} per trasmettere a più alta frequenza sono varie:
\begin{itemize}[label=\textcolor{PoliBlue}{\textbullet}]
    \item storicamente, le prime \textbf{trasmissioni transcontinentali} ebbero successo (pur non sapendone il motivo, almeno fino ad Appleton) perché questo tipo di onde elettromagnetiche si riflettono nell'\textbf{atmosfera} (ionosfera);
    \item per un motivo analogo, ma opposto, si comunica con lo \textbf{spazio} a frequenze ancora più alte, perché \textbf{l'atmosfera è trasparente a tali frequenze};
    \item di base ad alte frequenze possono essere veicolate più informazioni, dal momento che ci si possono permettere \textbf{larghezze di banda maggiori}, e quindi un \textbf{maggior data rate}.
\end{itemize}

\begin{teorema}{Teorema della codifica di canale (Shannon-Hartley)}{shannon}
Dal teorema della codifica di canale di Shannon, sappiamo che il \textbf{limite massimo della capacità di canale}, ovvero banalmente il tasso di trasferimento dati che possiamo raggiungere, è dato da:
\begin{equation}
    C = B \log_2 \left( 1 + \frac{S}{N} \right)
\end{equation}
\end{teorema}

Dunque, per incrementare la capacità $C$, possiamo \textbf{migliorare $S/N$}, e/oppure \textbf{allargare $B$} (ma in questo secondo caso attenzione al rumore, poiché allargando la banda di ricezione si cattura intrinsecamente una quantità maggiore di rumore termico).

Determinare una nomenclatura per i diversi gruppi di frequenze non è facile, dal momento che non c'è un approccio univoco: se ne può delineare una classificazione di massima.

\begin{itemize}[label=\textcolor{PoliBlue}{\textbullet}]
    \item Il blocco di frequenze fino a \qty{3}{\tera\hertz} (o \qty{0.1}{\milli\meter}) sono denominate \textbf{radiofrequenze} (RF): le onde radio però identificano le frequenze basse (sotto i \qty{300}{\mega\hertz}).
    \item Sopra i \unit{\giga\hertz} (sopra i \qty{300}{\mega\hertz} per l'esattezza) si parla di \textbf{microonde}.
\end{itemize}

La faccenda si complica perché a livello \textbf{internazionale} l'intero spettro è stato suddiviso in macrobande, la cui nomenclatura non è univoca.

\begin{figure}[H]
    \centering
    % Grafico logaritmico dello spettro (x = log10(frequenza in Hz))
    % Scala x: 1 unità = 1 decade (potenza di 10). 0 = 1 Hz, 13 = 10 THz
    \resizebox{\textwidth}{!}{
    \begin{tikzpicture}[font=\sffamily\scriptsize, xscale=1.1, yscale=0.8]
        
        % --- Griglia Logaritmica ---
        \foreach \d in {0,1,2,3,4,5,6,7,8,9,10,11,12,13} {
            % Linea principale (Decadi)
            \draw[gray!60, thin] (\d, -0.5) -- (\d, 5.5);
            % Linee secondarie (logaritmiche: 2, 3, 4... 9)
            \ifnum\d<13
                \foreach \m in {2,3,4,5,6,7,8,9} {
                    \pgfmathsetmacro{\x}{\d + log10(\m)}
                    \draw[gray!20, very thin] (\x, -0.2) -- (\x, 5.5);
                }
            \fi
        }
        
        % --- Etichette Asse X ---
        \node[below] at (0, -0.5) {1 Hz};
        \node[below] at (1, -0.5) {10 Hz};
        \node[below] at (2, -0.5) {100 Hz};
        \node[below] at (3, -0.5) {1 kHz};
        \node[below] at (4, -0.5) {10 kHz};
        \node[below] at (5, -0.5) {100 kHz};
        \node[below] at (6, -0.5) {1 MHz};
        \node[below] at (7, -0.5) {10 MHz};
        \node[below] at (8, -0.5) {100 MHz};
        \node[below] at (9, -0.5) {1 GHz};
        \node[below] at (10, -0.5) {10 GHz};
        \node[below] at (11, -0.5) {100 GHz};
        \node[below] at (12, -0.5) {1 THz};
        \node[below] at (13, -0.5) {10 THz};
        
        % --- Etichette Asse Y ---
        \node[left, font=\sffamily\bfseries\small] at (-0.2, 0.5) {ITU};
        \node[left, font=\sffamily\bfseries\small] at (-0.2, 2.5) {IEEE};
        \node[left, font=\sffamily\bfseries\small, align=right] at (-0.2, 4.5) {EU\\NATO\\US ECM};

        % ==========================================
        % 1. BANDA ITU (y tra 0 e 1)
        % Le frequenze ITU tagliano a 3*10^n. log10(3) = 0.477
        % ==========================================
        \def\yA{0} \def\yB{1}
        \filldraw[fill=ForestGreen!80!black, draw=white, opacity=0.8] (3.477, \yA) rectangle (4.477, \yB) node[midway, text=white] {VLF};
        \filldraw[fill=Purple!80!white, draw=white, opacity=0.8]      (4.477, \yA) rectangle (5.477, \yB) node[midway, text=white] {LF};
        \filldraw[fill=Cyan!70!blue, draw=white, opacity=0.8]         (5.477, \yA) rectangle (6.477, \yB) node[midway, text=white] {MF};
        \filldraw[fill=Orange!90!red, draw=white, opacity=0.8]        (6.477, \yA) rectangle (7.477, \yB) node[midway, text=white] {HF};
        \filldraw[fill=RoyalBlue!80, draw=white, opacity=0.8]         (7.477, \yA) rectangle (8.477, \yB) node[midway, text=white] {VHF};
        \filldraw[fill=Red!70, draw=white, opacity=0.8]               (8.477, \yA) rectangle (9.477, \yB) node[midway, text=white] {UHF};
        \filldraw[fill=ForestGreen!70, draw=white, opacity=0.8]       (9.477, \yA) rectangle (10.477, \yB) node[midway, text=white] {SHF};
        \filldraw[fill=Purple!70, draw=white, opacity=0.8]            (10.477, \yA) rectangle (11.477, \yB) node[midway, text=white] {EHF};
        \filldraw[fill=Cyan!70, draw=white, opacity=0.8]              (11.477, \yA) rectangle (12.477, \yB) node[midway, text=white] {THF};

        % ==========================================
        % 2. BANDA IEEE (y tra 2 e 3)
        % ==========================================
        \def\yC{2} \def\yD{3}
        \filldraw[fill=ForestGreen!80, draw=white, opacity=0.8] (6.477, \yC) rectangle (7.477, \yD) node[midway, text=black] {HF};
        \filldraw[fill=Purple!70!gray, draw=white, opacity=0.8] (7.477, \yC) rectangle (8.477, \yD) node[midway, text=white] {VHF};
        \filldraw[fill=Cyan!80!blue, draw=white, opacity=0.8]   (8.477, \yC) rectangle (9.000, \yD) node[midway, text=white] {UHF};
        \filldraw[fill=Orange, draw=white, opacity=0.8]         (9.000, \yC) rectangle (9.301, \yD) node[midway, text=black] {L};
        \filldraw[fill=RoyalBlue!70, draw=white, opacity=0.8]   (9.301, \yC) rectangle (9.602, \yD) node[midway, text=white] {S};
        \filldraw[fill=Red!70, draw=white, opacity=0.8]         (9.602, \yC) rectangle (9.903, \yD) node[midway, text=white] {C};
        \filldraw[fill=ForestGreen!70, draw=white, opacity=0.8] (9.903, \yC) rectangle (10.079, \yD) node[midway, text=black] {X};
        \filldraw[fill=Purple!70, draw=white, opacity=0.8]      (10.079, \yC) rectangle (10.255, \yD) node[midway, text=white] {Ku};
        \filldraw[fill=Cyan!70, draw=white, opacity=0.8]        (10.255, \yC) rectangle (10.423, \yD) node[midway, text=black] {K};
        \filldraw[fill=Orange!70, draw=white, opacity=0.8]      (10.423, \yC) rectangle (10.602, \yD) node[midway, text=black] {Ka};
        \filldraw[fill=RoyalBlue!40, draw=white, opacity=0.8]   (10.602, \yC) rectangle (10.875, \yD) node[midway, text=black] {V};
        \filldraw[fill=Red!40, draw=white, opacity=0.8]         (10.875, \yC) rectangle (11.041, \yD) node[midway, text=black] {W};
        \filldraw[fill=ForestGreen!30, draw=white, opacity=0.8] (11.041, \yC) rectangle (11.477, \yD) node[midway, text=black] {mm};

        % ==========================================
        % 3. BANDA EU/NATO/US ECM (y tra 4 e 5)
        % ==========================================
        \def\yE{4} \def\yF{5}
        \filldraw[fill=ForestGreen!70!black, draw=white, opacity=0.8] (0, \yE) rectangle (8.398, \yF) node[midway, text=black] {A}; % fino a 250MHz
        \filldraw[fill=Purple!70!gray, draw=white, opacity=0.8]       (8.398, \yE) rectangle (8.699, \yF) node[midway, text=white] {B}; % 500MHz
        \filldraw[fill=Cyan!80!blue, draw=white, opacity=0.8]         (8.699, \yE) rectangle (9.000, \yF) node[midway, text=white] {C}; % 1GHz
        \filldraw[fill=Orange, draw=white, opacity=0.8]               (9.000, \yE) rectangle (9.301, \yF) node[midway, text=black] {D}; % 2GHz
        \filldraw[fill=RoyalBlue!70, draw=white, opacity=0.8]         (9.301, \yE) rectangle (9.477, \yF) node[midway, text=white] {E}; % 3GHz
        \filldraw[fill=Red!70, draw=white, opacity=0.8]               (9.477, \yE) rectangle (9.602, \yF) node[midway, text=white] {F}; % 4GHz
        \filldraw[fill=ForestGreen!70, draw=white, opacity=0.8]       (9.602, \yE) rectangle (9.778, \yF) node[midway, text=black] {G}; % 6GHz
        \filldraw[fill=Purple!70, draw=white, opacity=0.8]            (9.778, \yE) rectangle (9.903, \yF) node[midway, text=white] {H}; % 8GHz
        \filldraw[fill=Cyan!70, draw=white, opacity=0.8]              (9.903, \yE) rectangle (10.000, \yF) node[midway, text=black] {I}; % 10GHz
        \filldraw[fill=Orange!70, draw=white, opacity=0.8]            (10.000, \yE) rectangle (10.301, \yF) node[midway, text=black] {J}; % 20GHz
        \filldraw[fill=RoyalBlue!40, draw=white, opacity=0.8]         (10.301, \yE) rectangle (10.602, \yF) node[midway, text=black] {K}; % 40GHz
        \filldraw[fill=Red!40, draw=white, opacity=0.8]               (10.602, \yE) rectangle (10.778, \yF) node[midway, text=black] {L}; % 60GHz
        \filldraw[fill=ForestGreen!30, draw=white, opacity=0.8]       (10.778, \yE) rectangle (11.000, \yF) node[midway, text=black] {M}; % 100GHz

    \end{tikzpicture}
    }
    \caption*{Classificazione dello spettro elettromagnetico secondo i principali standard internazionali.}
\end{figure}

\subsection*{Designazione delle Bande (Standard ITU)}

L'Unione Internazionale delle Telecomunicazioni (ITU) divide lo spettro radio in bande logaritmiche, molto usate per applicazioni broadcast, marittime e aeronautiche.

\begin{table}[H]
    \centering
    \small
    \renewcommand{\arraystretch}{1.3} % Dà un po' di respiro alle righe
    \begin{tabularx}{\textwidth}{>{\bfseries\color{PoliBlue}}c c c >{\raggedright\arraybackslash}X}
        \toprule
        \rowcolor{PoliBlue!10} 
        \textcolor{black}{Banda} & \textbf{Frequenze} & \textbf{Lunghezza d'onda} & \textbf{Uso Tipico} \\
        \midrule
        ELF & \qty{3}{\hertz} -- \qty{30}{\hertz} & \qty{100000}{\kilo\meter} -- \qty{10000}{\kilo\meter} & Trasmissione di potenza \\
        SLF & \qty{30}{\hertz} -- \qty{300}{\hertz} & \qty{10000}{\kilo\meter} -- \qty{1000}{\kilo\meter} & \\
        ULF & \qty{300}{\hertz} -- \qty{3000}{\hertz} & \qty{1000}{\kilo\meter} -- \qty{100}{\kilo\meter} & Frequenze vocali (Audio) \\
        VLF & \qty{3}{\kilo\hertz} -- \qty{30}{\kilo\hertz} & \qty{100}{\kilo\meter} -- \qty{10}{\kilo\meter} & Radionavigazione a largo raggio. Comunicazioni sottomarine. \\
        LF & \qty{30}{\kilo\hertz} -- \qty{300}{\kilo\hertz} & \qty{10}{\kilo\meter} -- \qty{1}{\kilo\meter} & Radiolocalizzazione marittima ed aeronautica. \\
        MF & \qty{300}{\kilo\hertz} -- \qty{3000}{\kilo\hertz} & \qty{1}{\kilo\meter} -- \qty{100}{\meter} & Comunicazioni aeree e marittime. Radionavigazione. Broadcast AM. \\
        HF & \qty{3}{\mega\hertz} -- \qty{30}{\mega\hertz} & \qty{100}{\meter} -- \qty{10}{\meter} & Collegamenti senza fili a lunga distanza fissi e mobili. Radioamatori (Onde Corte). \\
        VHF & \qty{30}{\mega\hertz} -- \qty{300}{\mega\hertz} & \qty{10}{\meter} -- \qty{1}{\meter} & Broadcast FM e TV. Collegamenti in visibilità. Radiomobili civili e militari. \\
        UHF & \qty{300}{\mega\hertz} -- \qty{3000}{\mega\hertz} & \qty{1}{\meter} -- \qty{10}{\centi\meter} & Ponti radio e radiomobili terrestri. Broadcast TV. Telefonia cellulare. Satelliti. \\
        SHF & \qty{3}{\giga\hertz} -- \qty{30}{\giga\hertz} & \qty{10}{\centi\meter} -- \qty{1}{\centi\meter} & Ponti radio terrestri. Satelliti. Radar. \\
        EHF & \qty{30}{\giga\hertz} -- \qty{300}{\giga\hertz} & \qty{1}{\centi\meter} -- \qty{1}{\milli\meter} & Onde millimetriche. Radar e sperimentazioni. \\
        \bottomrule
    \end{tabularx}
\end{table}

\subsection*{Bande ad Alta Frequenza (Radar, Spazio e Microonde)}

Salendo ancora più in alto con le frequenze (nel campo delle microonde), si entra in zone dello spettro fondamentali per le comunicazioni spaziali, i radar e i ponti radio ad altissima capacità. Per queste frequenze si adotta comunemente la letteratura IEEE.

\vspace{0.3cm}

\textbf{1. Frequenze per le comunicazioni spaziali (Satelliti e Deep Space)}
\begin{table}[H]
    \centering
    \small
    \renewcommand{\arraystretch}{1.2}
    \begin{tabular}{lcc}
        \toprule
        \rowcolor{PoliBlue!10}
        \textbf{Banda} & \multicolumn{2}{c}{\textbf{Frequenza (GHz)}} \\
        \rowcolor{PoliBlue!10}
        & \textbf{Downlink} & \textbf{Uplink} \\
        \midrule
        \textbf{S (deep space)} & $2.29 - 2.30$ & $2.11 - 2.12$ \\
        \textbf{S} & $2.20 - 2.29$ & $2.025 - 2.11$ \\
        \textbf{C} & $3.4 - 4.2$ & $5.925 - 6.425$ \\
        \textbf{X} & $7.25 - 7.75$ & $7.9 - 8.4$ \\
                   & $8.4 - 8.5$ & $7.145 - 7.235$ \\
        \textbf{Ku} & $10.7 - 12.75$ & $13.75 - 14.5$ \\
        \textbf{Ka} & $17.7 - 21.2$ & $27.5 - 31$ \\
                    & $25.5 - 27.5$ & \\
        \bottomrule
    \end{tabular}
\end{table}

\vspace{0.3cm}

\textbf{2. Classificazione IEEE completa (Microonde e Radar)}
\begin{table}[H]
    \centering
    \footnotesize
    \renewcommand{\arraystretch}{1.3}
    \begin{tabularx}{\textwidth}{>{\bfseries\color{PoliBlue}}c c c >{\raggedright\arraybackslash}X}
        \toprule
        \rowcolor{PoliBlue!10} 
        \textcolor{black}{Designazione} & \textbf{Gamma di Freq.} & \textbf{Lunghezza d'onda} & \textbf{Usi Tipici} \\
        \midrule
        L & \qty{1}{\giga\hertz} -- \qty{2}{\giga\hertz} & \qty{15}{\centi\meter} -- \qty{30}{\centi\meter} & Telemetria militare, GPS, telefonia mobile (GSM), radioamatori. \\
        S & \qty{2}{\giga\hertz} -- \qty{4}{\giga\hertz} & \qty{7.5}{\centi\meter} -- \qty{15}{\centi\meter} & Radar meteorologici e navali, satelliti, forni a microonde, Wi-Fi, Bluetooth, ZigBee. \\
        C & \qty{4}{\giga\hertz} -- \qty{8}{\giga\hertz} & \qty{3.75}{\centi\meter} -- \qty{7.5}{\centi\meter} & Telecomunicazioni radio a lunga distanza. \\
        X & \qty{8}{\giga\hertz} -- \qty{12}{\giga\hertz} & \qty{25}{\milli\meter} -- \qty{37.5}{\milli\meter} & Satelliti, radar, broadband terrestre, comunicazioni spaziali. \\
        Ku & \qty{12}{\giga\hertz} -- \qty{18}{\giga\hertz} & \qty{16.7}{\milli\meter} -- \qty{25}{\milli\meter} & Comunicazioni satellitari. \\
        K & \qty{18}{\giga\hertz} -- \qty{26.5}{\giga\hertz} & \qty{11.3}{\milli\meter} -- \qty{16.7}{\milli\meter} & Radar, comunicazioni satellitari, osservazioni astronomiche, radar automotive. \\
        Ka & \qty{26.5}{\giga\hertz} -- \qty{40}{\giga\hertz} & \qty{5.0}{\milli\meter} -- \qty{11.3}{\milli\meter} & Comunicazioni satellitari ad alta risoluzione. \\
        Q & \qty{33}{\giga\hertz} -- \qty{50}{\giga\hertz} & \qty{6.0}{\milli\meter} -- \qty{9.0}{\milli\meter} & Satelliti, comunicazioni terrestri a microonde, radioastronomia, radar automotive. \\
        U & \qty{40}{\giga\hertz} -- \qty{60}{\giga\hertz} & \qty{5.0}{\milli\meter} -- \qty{7.5}{\milli\meter} & \\
        V & \qty{50}{\giga\hertz} -- \qty{75}{\giga\hertz} & \qty{4.0}{\milli\meter} -- \qty{6.0}{\milli\meter} & Ricerca radar a onde millimetriche e spettroscopia. \\
        W & \qty{75}{\giga\hertz} -- \qty{110}{\giga\hertz} & \qty{2.7}{\milli\meter} -- \qty{4.0}{\milli\meter} & Satelliti, radar militari per tracciamento bersagli, radar automotive. \\
        F & \qty{90}{\giga\hertz} -- \qty{140}{\giga\hertz} & \qty{2.1}{\milli\meter} -- \qty{3.3}{\milli\meter} & Radioastronomia, radar avanzati, satelliti per broadcast televisivo (DBS). \\
        D & \qty{110}{\giga\hertz} -- \qty{170}{\giga\hertz} & \qty{1.8}{\milli\meter} -- \qty{2.7}{\milli\meter} & Radioastronomia, ponti radio ad altissima frequenza, remote sensing, armi a energia diretta. \\
        \bottomrule
    \end{tabularx}
\end{table}


\section{Radio e Frequenze}

Dunque, siccome le informazioni di interesse vengono veicolate tramite canali a \textbf{frequenza più alta rispetto alla banda base}, sono necessarie delle \textbf{conversioni} in frequenza.

\begin{itemize}[label=\textcolor{PoliBlue}{\textbullet}]
    \item Lo \textbf{scopo principale della radio} è proprio quello di permettere l'accesso alle trasmissioni attraverso onde elettromagnetiche a più alta frequenza.
\end{itemize}

Siccome nelle telecomunicazioni occorre compiere questo transito attraverso le alte frequenze \textbf{in due direzioni} (andata e ritorno), allora sono due i dispositivi d'interesse per una radio: \textbf{trasmettitore} e \textbf{ricevitore}.


\subsection*{Trasmettitore (Up-conversion)}

Un \textbf{trasmettitore (Tx o TX)} radio converte un segnale a bassa frequenza in un segnale ad alta frequenza (\textit{up-conversion}):
\begin{itemize}[label=\textcolor{PoliBlue}{\textbullet}]
    \item pensiamo al segnale corrispondente a una \textbf{voce}: un \textbf{microfono} trasduce il segnale sonoro in un segnale elettrico che passa alla radio, dove viene modificato in frequenza (e potenza) per finire all'antenna da cui avviene la \textbf{trasmissione};
    \item nell'esempio qui sotto si immagina una \textit{up-conversion} di un tono da una frequenza a una più alta.
\end{itemize}

\begin{figure}[H]
    \centering
    \begin{tikzpicture}[xscale=1.5, yscale=1.2, font=\sffamily]
        % Seed per rumore costante
        \pgfmathsetseed{111}
        
        % Assi
        \draw[-{Stealth}, ultra thick, red] (0,0) -- (0,2.5) node[left] {P};
        \draw[-{Stealth}, ultra thick, red] (0,0) -- (4.5,0) node[below right] {$\nu$};
        
        % Rumore di fondo
        \draw[thick, black!80] plot[domain=0.05:4.2, samples=150] (\x, {0.15 + 0.1*rand + 0.05*sin(\x*1000)});
        
        % Picco in Banda Base (Bassa Frequenza)
        \draw[thick, orange] (1, 0) -- (1, 1.5);
        \draw[thick, orange] (0.95, 0) -- (1, 1.5) -- (1.05, 0);
        
        % Picco in Radiofrequenza (Alta Frequenza)
        \draw[thick, orange] (3.5, 0) -- (3.5, 1.5);
        \draw[thick, orange] (3.45, 0) -- (3.5, 1.5) -- (3.55, 0);
        
        % Freccia Up-conversion
        \draw[-{Stealth[scale=1.5]}, line width=3pt, PoliBlue] (1.2, -0.4) to[bend right=20] (3.3, -0.4);
    \end{tikzpicture}
\end{figure}


\subsection*{Ricevitore (Down-conversion)}

Un \textbf{ricevitore (Rx o RX)} radio converte un segnale ad alta frequenza in un segnale a bassa frequenza (\textit{down-conversion}):
\begin{itemize}[label=\textcolor{PoliBlue}{\textbullet}]
    \item una antenna rileva un segnale che consegna alla radio: qui il segnale viene convertito a più bassa frequenza e passato a un trasduttore di onde elettriche in onde sonore, quale una \textbf{cassa acustica};
    \item nell'esempio qui sotto si immagina la \textit{down-conversion} di un tono a una frequenza più bassa.
\end{itemize}

\begin{figure}[H]
    \centering
    \begin{tikzpicture}[xscale=1.5, yscale=1.2, font=\sffamily]
        % Seed per rumore costante
        \pgfmathsetseed{222}
        
        % Assi
        \draw[-{Stealth}, ultra thick, red] (0,0) -- (0,2.5) node[left] {P};
        \draw[-{Stealth}, ultra thick, red] (0,0) -- (4.5,0) node[below right] {$\nu$};
        
        % Rumore di fondo
        \draw[thick, black!80] plot[domain=0.05:4.2, samples=150] (\x, {0.15 + 0.1*rand + 0.05*sin(\x*1200)});
        
        % Picco in Banda Base (Bassa Frequenza)
        \draw[thick, orange] (1, 0) -- (1, 1.5);
        \draw[thick, orange] (0.95, 0) -- (1, 1.5) -- (1.05, 0);
        
        % Picco in Radiofrequenza (Alta Frequenza)
        \draw[thick, orange] (3.5, 0) -- (3.5, 1.5);
        \draw[thick, orange] (3.45, 0) -- (3.5, 1.5) -- (3.55, 0);
        
        % Freccia Down-conversion
        \draw[-{Stealth[scale=1.5]}, line width=3pt, PoliBlue] (3.3, -0.4) to[bend left=20] (1.2, -0.4);
    \end{tikzpicture}
\end{figure}


\subsection*{Frequenza Intermedia (IF)}

Solitamente, soprattutto quando le trasmissioni avvengono alle alte frequenze sopra il GHz, il passaggio diretto da banda base a radiofrequenza (e viceversa) non avviene per limiti tecnologici nei dispositivi elettronici; bensì c'è un transito verso una \textbf{media frequenza}, detta appunto \textit{intermediate frequency} (\textbf{IF}).

Vengono mostrati un caso \textbf{Tx} (a sinistra) e un caso \textbf{Rx} (a destra).

\begin{figure}[H]
    \centering
    \begin{tikzpicture}[xscale=1.1, yscale=1.1, font=\sffamily\small]
        
        % Seed per il rumore
        \pgfmathsetseed{333}
        
        % =========================================
        % GRAFICO 1: Trasmettitore (Tx)
        % =========================================
        \begin{scope}[shift={(0,0)}]
            % Assi
            \draw[-{Stealth}, ultra thick, red] (0,0) -- (0,2.5) node[left] {P};
            \draw[-{Stealth}, ultra thick, red] (0,0) -- (4,0) node[below right] {$\nu$};
            
            % Rumore
            \draw[thick, black!80] plot[domain=0.05:3.8, samples=150] (\x, {0.12 + 0.08*rand});
            
            % Picco BB (Banda Base)
            \draw[thick, orange] (0.8, 0) -- (0.8, 1.5);
            \draw[thick, orange] (0.75, 0) -- (0.8, 1.5) -- (0.85, 0);
            \node[below, text=orange, font=\bfseries] at (0.8, -0.1) {BB};
            
            % Picco IF (Frequenza Intermedia)
            \draw[thick, ForestGreen] (2.0, 0) -- (2.0, 1.5);
            \draw[thick, ForestGreen] (1.95, 0) -- (2.0, 1.5) -- (2.05, 0);
            \node[below, text=ForestGreen, font=\bfseries] at (2.0, -0.1) {IF};
            
            % Picco RF (Radio Frequenza)
            \draw[thick, brown!80!black] (3.2, 0) -- (3.2, 1.5);
            \draw[thick, brown!80!black] (3.15, 0) -- (3.2, 1.5) -- (3.25, 0);
            \node[below, text=brown!80!black, font=\bfseries] at (3.2, -0.1) {RF};
            
            % Frecce Up-conversion a stadi
            \draw[-{Stealth[scale=1.2]}, line width=2pt, PoliBlue] (0.9, -0.6) to[bend right=30] (1.9, -0.6);
            \draw[-{Stealth[scale=1.2]}, line width=2pt, PoliBlue] (2.1, -0.6) to[bend right=30] (3.1, -0.6);
        \end{scope}

        % =========================================
        % GRAFICO 2: Ricevitore (Rx)
        % =========================================
        \begin{scope}[shift={(6,0)}]
            % Assi
            \draw[-{Stealth}, ultra thick, red] (0,0) -- (0,2.5) node[left] {P};
            \draw[-{Stealth}, ultra thick, red] (0,0) -- (4,0) node[below right] {$\nu$};
            
            % Rumore
            \draw[thick, black!80] plot[domain=0.05:3.8, samples=150] (\x, {0.12 + 0.08*rand});
            
            % Picco BB
            \draw[thick, orange] (0.8, 0) -- (0.8, 1.5);
            \draw[thick, orange] (0.75, 0) -- (0.8, 1.5) -- (0.85, 0);
            \node[below, text=orange, font=\bfseries] at (0.8, -0.1) {BB};
            
            % Picco IF
            \draw[thick, ForestGreen] (2.0, 0) -- (2.0, 1.5);
            \draw[thick, ForestGreen] (1.95, 0) -- (2.0, 1.5) -- (2.05, 0);
            \node[below, text=ForestGreen, font=\bfseries] at (2.0, -0.1) {IF};
            
            % Picco RF
            \draw[thick, brown!80!black] (3.2, 0) -- (3.2, 1.5);
            \draw[thick, brown!80!black] (3.15, 0) -- (3.2, 1.5) -- (3.25, 0);
            \node[below, text=brown!80!black, font=\bfseries] at (3.2, -0.1) {RF};
            
            % Frecce Down-conversion a stadi
            \draw[-{Stealth[scale=1.2]}, line width=2pt, PoliBlue] (3.1, -0.6) to[bend left=30] (2.1, -0.6);
            \draw[-{Stealth[scale=1.2]}, line width=2pt, PoliBlue] (1.9, -0.6) to[bend left=30] (0.9, -0.6);
        \end{scope}

    \end{tikzpicture}
\end{figure}


\subsection*{Filtraggio e Selezione della Banda}

Abbiamo già visto che in realtà la radio non converte un singolo tono, una singola frequenza, bensì una \textbf{banda di frequenze}.

\begin{itemize}[label=\textcolor{PoliBlue}{\textbullet}]
    \item Pensiamo ai segnali sonori: si tratta di una banda di frequenze e vogliamo trasmettere tutto ciò che è udibile; d'altra parte non vogliamo che si inseriscano altre frequenze possibile causa di rumore o interferenza.
    \item Perciò quando trasmettiamo o riceviamo un segnale di qualsiasi natura, vogliamo \textbf{selezionare le frequenze} che ci interessano ed escludere le altre: questa operazione si dice \textbf{filtraggio} e si opera tramite i filtri.
\end{itemize}

Quindi un \textbf{TX} e un \textbf{RX} selezionano una certa larghezza di banda del rispettivo segnale \textbf{in input}; altrettanto fanno \textbf{in output}, così da restituire solo quanto voluto (e così da massimizzare la potenza trasmessa\dots)

\begin{figure}[H]
    \centering
    \begin{tikzpicture}[xscale=1.1, yscale=1.1, font=\sffamily\small]
        
        \pgfmathsetseed{444}
        
        % =========================================
        % GRAFICO 1: Trasmettitore (Filtro Up)
        % =========================================
        \begin{scope}[shift={(0,0)}]
            % Assi
            \draw[-{Stealth}, ultra thick, red] (0,0) -- (0,2.5) node[left] {P};
            \draw[-{Stealth}, ultra thick, red] (0,0) -- (4.2,0) node[below right] {$\nu$};
            
            % Rumore
            \draw[thick, black!80] plot[domain=0.05:4.0, samples=150] (\x, {0.12 + 0.08*rand});
            
            % Picco 1
            \draw[thick, brown!80!black] (1, 0) -- (1, 1.3);
            \draw[thick, brown!80!black] (0.95, 0) -- (1, 1.3) -- (1.05, 0);
            
            % Filtro (Trapezio arancione) BB
            \draw[thick, orange] (0.4, 0) -- (0.5, 1.4) -- (1.5, 1.4) -- (1.6, 0);
            
            % Picco 2
            \draw[thick, brown!80!black] (3, 0) -- (3, 1.3);
            \draw[thick, brown!80!black] (2.95, 0) -- (3, 1.3) -- (3.05, 0);
            
            % Filtro (Trapezio arancione) RF
            \draw[thick, orange] (2.4, 0) -- (2.5, 1.4) -- (3.5, 1.4) -- (3.6, 0);
            
            % Freccia Up
            \draw[-{Stealth[scale=1.5]}, line width=3pt, PoliBlue] (1.2, -0.6) to[bend right=25] (2.8, -0.6);
        \end{scope}

        % =========================================
        % GRAFICO 2: Ricevitore (Filtro Down)
        % =========================================
        \begin{scope}[shift={(6.5,0)}]
            % Assi
            \draw[-{Stealth}, ultra thick, red] (0,0) -- (0,2.5) node[left] {P};
            \draw[-{Stealth}, ultra thick, red] (0,0) -- (4.2,0) node[below right] {$\nu$};
            
            % Rumore
            \draw[thick, black!80] plot[domain=0.05:4.0, samples=150] (\x, {0.12 + 0.08*rand});
            
            % Picco 1
            \draw[thick, brown!80!black] (1, 0) -- (1, 1.3);
            \draw[thick, brown!80!black] (0.95, 0) -- (1, 1.3) -- (1.05, 0);
            
            % Filtro BB
            \draw[thick, orange] (0.4, 0) -- (0.5, 1.4) -- (1.5, 1.4) -- (1.6, 0);
            
            % Picco 2
            \draw[thick, brown!80!black] (3, 0) -- (3, 1.3);
            \draw[thick, brown!80!black] (2.95, 0) -- (3, 1.3) -- (3.05, 0);
            
            % Filtro RF
            \draw[thick, orange] (2.4, 0) -- (2.5, 1.4) -- (3.5, 1.4) -- (3.6, 0);
            
            % Freccia Down
            \draw[-{Stealth[scale=1.5]}, line width=3pt, PoliBlue] (2.8, -0.6) to[bend left=25] (1.2, -0.6);
        \end{scope}

    \end{tikzpicture}
\end{figure}

\subsection*{Radio e Livelli di Potenza}

Fin dall'inizio si è detto che la radio agisce, oltre che sulla frequenza, anche sulla \textbf{potenza}. 
Per esempio, quando parliamo regoliamo il volume della voce in base a due parametri:
\begin{itemize}[label=\textcolor{PoliBlue}{\textbullet}]
    \item la \textbf{distanza} del ricevente: maggiore questa distanza, più abbiamo bisogno di alzare il volume, cioè la potenza del segnale;
    \item il \textbf{livello di rumore}: maggiore è il brusio, più vogliamo smarcarci dal rumore di fondo. Viceversa, chi ascolta vuole distinguere una voce dal rumore di fondo.
\end{itemize}

Altrettanto fa la radio: infatti, siccome la distanza e il rumore lungo il tragitto Tx-to-Rx alterano il segnale:
\begin{itemize}[label=\textcolor{PoliBlue}{\textbullet}]
    \item il \textbf{trasmettitore} cerca di \textbf{alzare} la potenza trasmessa il più possibile senza alterare il segnale;
    \item il \textbf{ricevitore} cerca di \textbf{recuperare} il segnale (ormai attenuato dalla distanza) e \textbf{alzarne} la potenza rispetto al rumore.
\end{itemize}

\begin{figure}[H]
    \centering
    \begin{tikzpicture}[xscale=1.1, yscale=1.1, font=\sffamily\small]
        
        \pgfmathsetseed{555}
        
        % =========================================
        % GRAFICO 1: Trasmettitore (Amplificazione + Up)
        % =========================================
        \begin{scope}[shift={(0,0)}]
            % Assi
            \draw[-{Stealth}, ultra thick, red] (0,0) -- (0,3) node[left] {P};
            \draw[-{Stealth}, ultra thick, red] (0,0) -- (4.2,0) node[below right] {$\nu$};
            
            % Rumore
            \draw[thick, black!80] plot[domain=0.05:4.0, samples=150] (\x, {0.12 + 0.08*rand});
            
            % Segnale BB (Bassa Potenza)
            \draw[thick, brown!80!black] (1, 0) -- (1, 0.8);
            \draw[thick, brown!80!black] (0.95, 0) -- (1, 0.8) -- (1.05, 0);
            \draw[thick, orange] (0.4, 0) -- (0.5, 0.9) -- (1.5, 0.9) -- (1.6, 0); % Filtro
            
            % Segnale RF (Alta Potenza - Amplificato)
            \draw[thick, brown!80!black] (3, 0) -- (3, 2.5);
            \draw[thick, brown!80!black] (2.95, 0) -- (3, 2.5) -- (3.05, 0);
            \draw[thick, orange] (2.4, 0) -- (2.5, 2.6) -- (3.5, 2.6) -- (3.6, 0); % Filtro
            
            % Freccia Up
            \draw[-{Stealth[scale=1.5]}, line width=3pt, PoliBlue] (1.2, -0.6) to[bend right=25] (2.8, -0.6);
        \end{scope}

        % =========================================
        % GRAFICO 2: Ricevitore (Amplificazione + Down)
        % =========================================
        \begin{scope}[shift={(6.5,0)}]
            % Assi
            \draw[-{Stealth}, ultra thick, red] (0,0) -- (0,3) node[left] {P};
            \draw[-{Stealth}, ultra thick, red] (0,0) -- (4.2,0) node[below right] {$\nu$};
            
            % Rumore
            \draw[thick, black!80] plot[domain=0.05:4.0, samples=150] (\x, {0.12 + 0.08*rand});
            
            % Segnale RF Ricevuto (Attenuato, Bassa Potenza)
            \draw[thick, brown!80!black] (3, 0) -- (3, 0.7);
            \draw[thick, brown!80!black] (2.95, 0) -- (3, 0.7) -- (3.05, 0);
            \draw[thick, orange] (2.4, 0) -- (2.5, 0.8) -- (3.5, 0.8) -- (3.6, 0); % Filtro
            
            % Segnale BB Recuperato (Amplificato, Alta Potenza)
            \draw[thick, brown!80!black] (1, 0) -- (1, 1.8);
            \draw[thick, brown!80!black] (0.95, 0) -- (1, 1.8) -- (1.05, 0);
            \draw[thick, orange] (0.4, 0) -- (0.5, 1.9) -- (1.5, 1.9) -- (1.6, 0); % Filtro
            
            % Freccia Down
            \draw[-{Stealth[scale=1.5]}, line width=3pt, PoliBlue] (2.8, -0.6) to[bend left=25] (1.2, -0.6);
        \end{scope}

    \end{tikzpicture}
\end{figure}


% ==================================================
% NUOVA SEZIONE: ARCHITETTURE DEI RICEVITORI
% ==================================================
\section{Architetture dei Ricevitori Radio}

\subsection{Ricevitore Eterodina (Conversione Diretta)}

La denominazione \textbf{eterodina} fu inventata all'inizio del 1900 dal canadese Reginald Fessenden, con derivazione dal greco: \textit{heteros} = differente e \textit{dynamis} = potenza.

Questo termine si riferisce a un ricevitore per segnali emessi da trasmettitori ad arco sintonizzati:
\begin{itemize}[label=\textcolor{PoliBlue}{\textbullet}]
    \item il ricevitore era composto da un \textbf{mixer non lineare} e da un \textbf{oscillatore locale} costituito da un alternatore in alta frequenza (massimo \qty{100}{\kilo\hertz});
    \item fu il primo ricevitore a \textbf{conversione diretta}: la conversione arrivava direttamente in bassa frequenza, cioè si ascoltava direttamente con la cuffia il segnale telegrafico (Morse) ricevuto;
    \item tuttavia, l'alternatore usato come oscillatore locale si comportava a sua volta come un trasmettitore, ovvero \textbf{disturbava} il segnale d'interesse: il sistema fu abbandonato per passare alla supereterodina.
\end{itemize}

\begin{figure}[H]
    \centering
    \begin{tikzpicture}[font=\sffamily\small, >=Stealth, thick]
        
        % Nodo 1: Antenna (simbolo manuale)
        \coordinate (ant_base) at (0, 0);
        \draw (ant_base) -- (0, 1.5);
        \draw (0, 1.5) -- (-0.3, 1.8);
        \draw (0, 1.5) -- (0, 1.9);
        \draw (0, 1.5) -- (0.3, 1.8);
        \coordinate (ant_out) at (0, 0.5);
        
        % Nodo 2: Mixer
        \node[circle, draw, minimum size=1.2cm, fill=white] (mixer) at (2, 0.5) {};
        \draw (mixer.north west) -- (mixer.south east);
        \draw (mixer.north east) -- (mixer.south west);
        \node[above=0.1cm of mixer, font=\tiny] {Dispositivo non lineare};
        
        % Nodo 3: Oscillatore Locale
        \node[circle, draw, minimum size=1.2cm, fill=white] (osc) at (2, -2) {};
        \draw (1.7, -2) sin (1.85, -1.8) cos (2, -2) sin (2.15, -2.2) cos (2.3, -2);
        \draw[->] (1.5, -2.5) -- (2.5, -1.5); % Freccia di sintonizzazione
        
        % Nodo 4: Audio (Altoparlante)
        \coordinate (speaker_in) at (4.5, 0.5);
        \draw[fill=gray!20] (4.5, 0.2) rectangle (4.8, 0.8);
        \draw[fill=gray!20] (4.8, 0.2) -- (5.3, -0.2) -- (5.3, 1.2) -- (4.8, 0.8);
        \node[above=0.4cm] at (5, 0.5) {\scriptsize AUDIO};
        
        % Collegamenti
        \draw[->] (ant_out) -- (mixer.west);
        \draw[->] (osc.north) -- (mixer.south);
        \draw[->] (mixer.east) -- (speaker_in);
        
        % Etichette con box tratteggiato
        \node[draw, dotted, fill=white, inner sep=4pt] (label_ant) at (-1.5, 1.5) {\scriptsize Antenna};
        \draw[dotted] (label_ant.east) -- (-0.1, 1.5);
        
        \node[draw, dotted, fill=white, inner sep=4pt] (label_osc) at (2, -3.5) {\scriptsize Oscillatore Locale};
        \draw[dotted] (label_osc.north) -- (osc.south);
        
    \end{tikzpicture}
    \caption*{Schema a blocchi semplificato di un ricevitore eterodina a conversione diretta.}
\end{figure}


\subsection{Ricevitore Supereterodina}

Negli anni 1910/1915 Edwin Armstrong e Levy indipendentemente modificarono il principio della eterodina, costruendo l'oscillatore locale con un circuito a triodo (inventato da de Forest nel 1913). 

Invece di convertire la frequenza radio in ingresso dall'antenna direttamente in bassa frequenza, si fece una \textbf{conversione alla frequenza intermedia (IF)} fissa di \qty{50}{\kilo\hertz}. 
A questa frequenza IF si applicò un \textbf{amplificatore selettivo}, composto da filtro e amplificatore, prima che il trasduttore/detector ne estraessero il segnale acustico da inviare agli altoparlanti. Il sistema fu chiamato \textbf{supereterodina} (super inteso per supersonico).

\begin{nota}
    Il grandissimo vantaggio di questa tecnica consiste nel fatto che le diverse frequenze provenienti dalle diverse stazioni trasmittenti vengono \textbf{tutte convertite alla stessa frequenza intermedia} prima dell'amplificazione e del filtraggio. Gli stadi di amplificazione e filtraggio sono quindi più semplici ed economici da produrre perché operano a una frequenza intermedia fissa: non hanno più bisogno di essere regolabili per adattarsi alle diverse frequenze delle stazioni trasmittenti.
\end{nota}

In pratica il funzionamento è il seguente:
\begin{itemize}[label=\textcolor{PoliBlue}{\textbullet}]
    \item il segnale ricevuto (\(f_{RF}\)) viene combinato nel mixer con il segnale generato da un oscillatore locale (\(f_{LO}\)) e convertito a un nuovo valore di frequenza. 
    \item Questa nuova frequenza intermedia è tipicamente la \textbf{"differenza"} fra la frequenza dell'oscillatore locale e quella del segnale ricevuto:
    \begin{equation}
        f_{IF} = | f_{RF} - f_{LO} |
    \end{equation}
    \item \textcolor{gray}{\textit{Nota: matematicamente è presente anche la frequenza "somma" ($f_{RF} + f_{LO}$), che talvolta viene utilizzata in specifiche architetture al posto della differenza.}}
\end{itemize}


\subsection{Schema a Blocchi di un Ricevitore Supereterodina}

Per comprendere meglio il funzionamento di un ricevitore radio moderno, analizziamo il suo schema a blocchi completo. I concetti fondamentali alla base dell'architettura sono:
\begin{itemize}[label=\textcolor{PoliBlue}{\textbullet}]
    \item il segnale in ingresso dall'antenna si trova a una frequenza \textbf{RF > IF};
    \item al centro della catena, si converte la frequenza tramite l'oscillatore locale (LO) facendola passare da RF a IF;
    \item l'uscita IF deve avere un alto \textbf{SNR} (rapporto segnale/rumore) per essere inviata all'utilizzatore/trasduttore finale.
\end{itemize}

\begin{figure}[H]
    \centering
    \resizebox{\textwidth}{!}{%
    \begin{tikzpicture}[font=\sffamily\scriptsize, >=Stealth]
        
        % ==========================================
        % MACRO PER DISEGNARE I COMPONENTI
        % ==========================================
        % Amplificatore (Triangolo orizzontale)
        \newcommand{\drawamp}[3]{
            \draw[fill=white, thick] (#1, #2-0.6) -- (#1+1.2, #2) -- (#1, #2+0.6) -- cycle;
            \node at (#1+0.4, #2) {\textbf{#3}};
        }
        % Filtro (Rettangolo con onde)
        \newcommand{\drawfilter}[3]{
            \draw[fill=white, thick] (#1-0.6, #2-0.6) rectangle (#1+0.6, #2+0.6);
            \node[above] at (#1, #2+0.6) {\textbf{#3}};
            % Onde del filtro
            \draw[thick] (#1-0.4, #2+0.2) sin (#1-0.2, #2+0.35) cos (#1, #2+0.2) sin (#1+0.2, #2+0.05) cos (#1+0.4, #2+0.2);
            \draw[thick] (#1-0.4, #2-0.1) sin (#1-0.2, #2+0.05) cos (#1, #2-0.1) sin (#1+0.2, #2-0.25) cos (#1+0.4, #2-0.1);
            \draw[thick] (#1-0.4, #2-0.4) sin (#1-0.2, #2-0.25) cos (#1, #2-0.4) sin (#1+0.2, #2-0.55) cos (#1+0.4, #2-0.4);
            % Linea sbarrata (simbolo passa-banda)
            \draw[thick] (#1-0.3, #2+0.5) -- (#1+0.3, #2-0.1);
        }
        
        % ==========================================
        % ZONE COLORATE (Evidenziazioni stile slide)
        % ==========================================
        % 1. Front-End (Verde)
        \draw[ForestGreen, ultra thick, dashed, rounded corners] (0.8, -1.8) rectangle (6.5, 2.0);
        \node[ForestGreen, font=\sffamily\bfseries\large, anchor=south west] at (0.8, 2.0) {1. RX Front-End};

        % 2. Down-Conversion (Rosso)
        \draw[OpticRed, ultra thick, dashed, rounded corners] (6.8, -5.5) rectangle (8.8, 1.5);
        \node[OpticRed, font=\sffamily\bfseries\large, anchor=south] at (7.8, 1.5) {2. Down-Conversion};

        % 3. IF Section (Blu)
        \draw[PoliBlue, ultra thick, dashed, rounded corners] (9.2, -1.8) rectangle (13.5, 1.5);
        \node[PoliBlue, font=\sffamily\bfseries\large, anchor=south west] at (9.2, 1.5) {3. RX IF};


        % ==========================================
        % CATENA PRINCIPALE (Da sinistra a destra)
        % ==========================================
        % Input RF
        \node[left, text=OpticRed, font=\sffamily\bfseries\large] at (0, 1.5) {RF input};
        \filldraw[black] (0, 0) circle (2pt);
        \draw[thick, ->] (0, 0) -- (1, 0);
        
        % Testi Input
        \node[above, text=OpticRed, align=center] at (0, 0.2) {-115 dBm\\to\\-80 dBm};
        \node[below, align=center] at (0, -0.2) {\textbf{RX}\\27.5 - 30 GHz\\[0.2cm] CH1 - 27.7 GHz\\CH12 - 29.9 GHz\\Ad. Ch - 100 MHz};

        % LNA 1
        \drawamp{1}{0}{LNA 1}
        \draw[thick, ->] (2.2, 0) -- (3, 0);
        \node[below, text=OpticRed] at (1.6, -0.8) {25 dB};

        % LNA 2
        \drawamp{3}{0}{LNA 2}
        \draw[thick, ->] (4.2, 0) -- (5.2, 0);
        \node[below, text=OpticRed] at (3.6, -0.8) {25 dB};

        % Filtro F1
        \drawfilter{5.8}{0}{F1}
        \draw[thick, ->] (6.4, 0) -- (7.3, 0);
        \node[below, text=OpticRed, align=center] at (5.8, -0.8) {\textcolor{purple}{2500 MHz BW}\\-3 dB};

        % Mixer M1
        \node[circle, draw, thick, fill=white, minimum size=1cm] (mixer) at (7.8, 0) {};
        \draw[thick] (mixer.north west) -- (mixer.south east);
        \draw[thick] (mixer.north east) -- (mixer.south west);
        \node[above] at (7.8, 0.6) {\textbf{M1}};
        \draw[thick, ->] (8.3, 0) -- (9.4, 0);
        \node[above, font=\tiny\bfseries] at (8.85, 0) {IF1 - TBD};

        % Filtro F2
        \drawfilter{10.0}{0}{F2}
        \draw[thick, ->] (10.6, 0) -- (11.5, 0);
        \node[below, text=OpticRed, align=center] at (10.0, -0.8) {\textcolor{black}{TBD}\\\textcolor{purple}{200 MHz BW}\\-2 dB};

        % AMP 1
        \drawamp{11.5}{0}{AMP 1}
        \node[below, text=OpticRed] at (12.1, -0.8) {16 dB};

        % Output IF
        \draw[thick] (12.7, 0) -- (13.2, 0) -- (13.2, -1.5);
        \draw[thick, ->] (13.2, -1.5) -- (14.2, -1.5);
        \node[above, text=OpticRed, align=center] at (13.7, -1.4) {-64 dBm\\to\\-29 dBm};
        \node[right, text=OpticRed, font=\sffamily\bfseries\large] at (14.2, -2.5) {IF output};


        % ==========================================
        % CATENA OSCILLATORE LOCALE (LO) - Dal basso
        % ==========================================
        % Collegamento al Mixer
        \draw[thick, <-] (7.8, -0.5) -- (7.8, -1.0);
        
        % Moltiplicatore x2 (Triangolo verso l'alto)
        \draw[fill=white, thick] (7.4, -1.8) -- (8.2, -1.8) -- (7.8, -1.0) -- cycle;
        \node at (7.8, -1.5) {\textbf{x2}};
        \draw[thick, <-] (7.8, -1.8) -- (7.8, -2.2);
        
        % Buffer / Amplificatore LO
        \draw[fill=white, thick] (7.5, -2.8) -- (8.1, -2.8) -- (7.8, -2.2) -- cycle;
        \draw[thick, <-] (7.8, -2.8) -- (7.8, -3.3);
        \node[left, text=OpticRed] at (7.5, -2.5) {-10 dB};
        \node[right, font=\tiny\bfseries] at (7.9, -2.9) {TBD - TBD GHz};

        % Blocco LO (VCO + PLL)
        \draw[dashed, thick, fill=gray!5] (6.9, -5.0) rectangle (8.9, -3.3);
        \node[above left] at (7.5, -3.3) {\textbf{LO1}};
        
        % VCO
        \node[circle, draw, thick, fill=white, minimum size=0.8cm] (vco) at (7.8, -3.7) {};
        \draw[thick] (7.5, -3.7) sin (7.65, -3.55) cos (7.8, -3.7) sin (7.95, -3.85) cos (8.1, -3.7);
        \node[above] at (7.8, -3.3) {\textbf{VCO}};
        \draw[thick, <-] (7.8, -4.1) -- (7.8, -4.4);
        
        % Loop VCO -> PLL
        \draw[thick] (7.8, -3.1) -- (7.1, -3.1) -- (7.1, -4.6) -- (7.3, -4.6);
        
        % PLL
        \draw[fill=white, thick] (7.3, -4.8) rectangle (8.3, -4.3);
        \node at (7.8, -4.55) {\textbf{PLL}};
        
        % f_ref
        \draw[thick, OpticRed, <-] (7.8, -5.0) -- (7.8, -5.4) node[right, text=black] {$f_{\text{ref}}$};

    \end{tikzpicture}
    }
\end{figure}

Il diagramma soprastante può essere idealmente suddiviso in tre macro-sezioni funzionali, di cui andiamo ora ad analizzare il comportamento.

\subsubsection*{1. RX Front-End}
La primissima parte del circuito, posta subito a valle dell'antenna, prende il nome di \textit{Front-End}. È composto tipicamente da:
\begin{itemize}[label=\textcolor{ForestGreen}{\textbullet}]
    \item \textbf{Amplificatori a basso rumore (LNA - Low Noise Amplifier)}: sono amplificatori speciali capaci di elevare il livello di potenza di un segnale estremamente debole (fino a \qty{-115}{\deci\bel m}) introdotto dall'antenna, cercando di \textbf{non alzare il rumore} termico di fondo.
    \item \textbf{Filtro reiezione di immagine (F1)}: un filtro passa-banda che ripulisce il segnale in ingresso da interferenze fuori banda prima della conversione.
\end{itemize}
\begin{nota}[colback=ForestGreen!5, colframe=ForestGreen, title={\textbf{Cosa succede nel Front-End?}}]
    Qui cambia in modo significativo la \textbf{potenza} (viene aumentata fortemente) e viene in parte \textbf{selezionata la banda} di frequenza desiderata. Tuttavia, \textbf{non cambia la frequenza} portante del segnale.
\end{nota}

\subsubsection*{2. Down-Conversion}
Il cuore del sistema supereterodina è la conversione di frequenza. I componenti chiave sono:
\begin{itemize}[label=\textcolor{OpticRed}{\textbullet}]
    \item \textbf{Mixer (Mescolatore M1)}: è un circuito intrinsecamente non lineare che opera matematicamente una \textbf{moltiplicazione} tra segnali. Nel dominio delle frequenze, questo si traduce in somme o sottrazioni di frequenze.
    \item \textbf{Oscillatore Locale (LO)}: deve generare e fornire un segnale stabile e pulito al mixer per convertire il segnale ricevuto (a frequenza RF) nel nuovo segnale a frequenza IF. Solitamente, per garantire assoluta precisione, questo segnale è comandato e stabilizzato da un \textbf{PLL} (\textit{Phase-Locked Loop}) accoppiato a un \textbf{VCO} (\textit{Voltage-Controlled Oscillator}).
\end{itemize}
\begin{nota}[colback=OpticRed!5, colframe=OpticRed, title={\textbf{Cosa succede nella Down-Conversion?}}]
    Questo è l'unico stadio in cui \textbf{cambia la frequenza} del segnale (da RF a IF). La potenza subisce variazioni modeste, solitamente una leggera perdita introdotta dal processo di miscelazione (\textit{conversion loss}).
\end{nota}

\subsubsection*{3. RX IF (Frequenza Intermedia)}
L'ultima fase del ricevitore opera su una frequenza più bassa e fissa (la IF). È composto da:
\begin{itemize}[label=\textcolor{PoliBlue}{\textbullet}]
    \item \textbf{Filtro selettore (F2)}: un filtro passa-banda molto stretto e preciso (essendo la IF fissa, questo filtro può essere ottimizzato per prestazioni eccellenti).
    \item \textbf{Amplificatore/i (AMP 1)}: stadi di amplificazione convenzionali che portano il segnale al livello necessario per il detector o il convertitore analogico-digitale (ADC).
\end{itemize}
\begin{nota}[colback=PoliBlue!5, colframe=PoliBlue, title={\textbf{Cosa succede nello stadio IF?}}]
    Viene \textbf{selezionata finemente la banda} di frequenza isolando esclusivamente il canale di interesse, e viene \textbf{aumentata la potenza} di quanto necessario all'utilizzatore. Anche qui, \textbf{non c'è spostamento di frequenza}.
\end{nota}

\end{document}